\SourcePage{371}
\SGAsection{6}{子函子与群子概形}
\footnote{关于同一主题，另见第 XI 讲第 6 节的结果；它们本来宜置于本讲 VIB 中，其中特别给出了给定\Term{group-scheme}{群概形}的某些\Term{group-subfunctor}{群子函子}的\Term{representability}{可表性}判据。编者注（48）：这些结果已插入下文（6.5.2 至 6.5.5）。}
\begin{SGAresult}{定义}{6.1}
(i) 设 $X$ 为一个 $S$-\Term{functor}{函子}（即从 $(\mathrm{Sch}/S)^\circ$ 到 $(\mathrm{Ens})$ 的\Term{functor}{函子}），$G$ 为一个 $S$-\Term{group-functor}{群函子}，$u,v$ 为从 $X$ 到 $G$ 的两个 $S$-\Term{morphism}{态射}。称以下方式定义的 $G$ 的 $S$-\Term{subfunctor}{子函子}为 $u$ 到 $v$ 的\Term{transporter}{输送子}，记作 $\operatorname{Transp}(u,v)$：
\SourcePage{372}
\[
\begin{aligned}
{}&\operatorname{Transp}(u,v)(S')\\
{}&=\{g\in G(S')\mid (\operatorname{int}g)\circ u_{S'}=v_{S'}\}\\
{}&=\{g\in G(S')\mid g_{S''}u_{S''}(x)g_{S''}^{-1}=v_{S''}(x),\\
{}&\hspace{39mm}\forall x\in X(S''),\ S''\to S'\}.
\end{aligned}
\]
特别地，$\operatorname{Transp}(u,u)$ 是 $G$ 的一个 $S$-\Term{group-subfunctor}{群子函子}，称为 $u$ 的\Term{centralizer}{中心化子}，记作 $\operatorname{Centr}(u)$。

(ii) 设 $G$ 为一个 $S$-\Term{group-functor}{群函子}，$X,Y$ 为 $G$ 的两个 $S$-\Term{subfunctor}{子函子}；称下面定义的 $G$ 的 $S$-\Term{subfunctor}{子函子}为 $X$ 到 $Y$ 的\Term{transporter}{输送子}（相应地，\Term{strict-transporter}{严格输送子}），记作 $\operatorname{Transp}_G(X,Y)$（相应地，$\operatorname{Transp}^{\mathrm{str}}_G(X,Y)$）：
\[
\begin{aligned}
\operatorname{Transp}_G(X,Y)(S')
 &=\{g\in G(S')\mid(\operatorname{int}g)(X_{S'})\subset Y_{S'}\}\\
 &=\{g\in G(S')\mid g_{S''}X(S'')g_{S''}^{-1}\subset Y(S''),\ \forall S''\to S'\},\\
\operatorname{Transp}^{\mathrm{str}}_G(X,Y)(S')
 &=\{g\in G(S')\mid(\operatorname{int}g)(X_{S'})=Y_{S'}\}\\
 &=\{g\in G(S')\mid g_{S''}X(S'')g_{S''}^{-1}=Y(S''),\ \forall S''\to S'\}.
\end{aligned}
\]
注意有
\[
\operatorname{Transp}^{\mathrm{str}}_G(X,Y)
 =\operatorname{Transp}_G(X,Y)\cap c\bigl(\operatorname{Transp}_G(Y,X)\bigr),
\]
其中 $c$ 表示 $G$ 的\Term{inversion-morphism}{求逆态射}。\footnote{编者注（49）：若 $u:X\to G$ 与 $v:Y\to G$ 是任意的 $S$-\Term{functor}{函子}\Term{morphism}{态射}，令 $u(X)$ 为 $u$ 的\Term{image-functor}{像函子}，定义为 $u(X)(S')=u(S')(X(S'))\subset G(S')$。它是 $G$ 的一个\Term{subfunctor}{子函子}，$v(Y)$ 也一样。因此可以考虑 $u$ 的像到 $v$ 的像的\Term{transporter}{输送子} $\operatorname{Transp}_G(u(X),v(Y))$。由此可见，定义 (ii) 不必局限于\Term{subfunctor}{子函子}，即不必要求 $u,v$ 是\Term{monomorphism}{单态射}。这一\Term{restriction}{限制}使原文的假设有时出现重复，例如：“设 $u,v:X\to G$ 为 $S$-\Term{functor}{函子}\Term{morphism}{态射}，$w:H\to G$ 与 $w':K\to G$ 为\Term{monomorphism}{单态射}，则 $\operatorname{Transp}(u,v)$、$\operatorname{Centr}_G(w)=\operatorname{Centr}_G H$ 与 $\operatorname{Transp}_G(H,K)$ 满足……”。改用 $\operatorname{Transp}_G(u(X),v(Y))$ 即可避免这些重复。我们在 6.2 及后面的 10.11 中作了这样的修改。}

(iii) 设 $G$ 为一个 $S$-\Term{group-functor}{群函子}，$H$ 为 $G$ 的一个 $S$-\Term{subfunctor}{子函子}，$i:H\to G$ 为\Term{canonical}{典范} $S$-\Term{morphism}{态射}。称下列 $S$-\Term{group-subfunctor}{群子函子}为 $H$ 在 $G$ 中的\Term{centralizer}{中心化子}与\Term{normalizer}{正规化子}：
\[
\operatorname{Centr}_G H=\operatorname{Centr}(i)=\operatorname{Transp}(i,i),
\qquad \operatorname{Norm}_G H=\operatorname{Transp}^{\mathrm{str}}_G(H,H).
\]
最后，称 $S$-\Term{group-functor}{群函子} $\operatorname{Centr}(\operatorname{id}_G)=\operatorname{Centr}_G G$ 为 $G$ 的\Term{center}{中心}，记作 $\operatorname{Centr}G$。
\end{SGAresult}

\textbf{注记 6.1.1。}\footnote{编者注（50）：本注记为新增，命题 6.2 中隐含地用到了它。}
由定义可知，\Term{functor}{函子} $\operatorname{Transp}(u,v)$ 与 $\operatorname{Transp}_G(X,Y)$（因而也包括 $\operatorname{Transp}^{\mathrm{str}}_G(X,Y)$）与\Term{base-change}{基变换}交换：
\SourcePage{373}
对任意 $S'\to S$，若 $G',X',Y',u',v'$ 由 $G,X,Y,u,v$ 经\Term{base-change}{基变换}得到，则
\[
\operatorname{Transp}(u,v)_{S'}=\operatorname{Transp}(u',v'),
\qquad \operatorname{Transp}(X,Y)_{S'}=\operatorname{Transp}(X',Y').
\]

\begin{SGAresult}{命题}{6.2}
\footnote{编者注（51）：如编者注（49）所述，我们修改了本命题的表述。}
设 $G$ 为一个 $S$-\Term{group}{群}。对\Term{functor}{函子} $G$ 的一个\Term{subfunctor}{子函子} $T$，考虑以下性质：
\[
(+_f)\qquad \text{对任意 }s\in S,\ T_s\text{ 可由 }G_s\text{ 的一个闭子概形表示。}
\]
设 $u,w:X\to G$ 与 $v:Y\to G$ 为 $S$-\Term{scheme}{概形}\Term{morphism}{态射}。则：

(i) $\operatorname{Transp}(u,w)$ 与 $\operatorname{Centr}(u)=\operatorname{Transp}(u,u)$ 满足条件 $(+_f)$。

(ii) 若对任意 $s\in S$，$v_s$ 都是\Term{closed-immersion}{闭浸入}，则 $\operatorname{Transp}_G(u(X),v(Y))$ 与 $\operatorname{Norm}_G v(Y)$ 满足条件 $(+_f)$。

(iii) 若对任意 $s\in S$，$u_s,v_s$ 都是\Term{closed-immersion}{闭浸入}，则 $\operatorname{Transp}^{\mathrm{str}}_G(X,Y)$ 满足条件 $(+_f)$。
\end{SGAresult}
这由注记 6.1.1 和下面的推论 6.2.5 得出。\footnote{编者注（52）：原文引用了第 VIII 讲第 6 节的结果。为方便读者，我们在下面的 6.2.1 至 6.2.5 中重录了这些结果（VIII，6.3 和 6.8 除外）。事实上，A. Grothendieck 已在 VIII.6 开头的注中建议：“本节不依赖于\Term{diagonalizable-group}{可对角化群}的理论；它的自然位置应在 VIB 中。”}

\begin{SGAresult}{定义}{6.2.1}
设 $f:X\to S$ 为一个\Term{scheme}{概形}\Term{morphism}{态射}。若存在 $S$ 的一个\Term{affine-open}{仿射开集}\Term{covering}{覆盖} $(S_i)$，对每个 $i$ 存在一个在 $S_i$ 上\Term{affine}{仿射}且\Term{faithfully-flat}{忠实平坦}的 $S_i$-\Term{scheme}{概形} $S'_i$，以及 $X'_i=X\times_S S'_i$ 的一个\Term{affine-open}{仿射开集}\Term{covering}{覆盖} $(X'_{ij})_j$，使每个 $X'_{ij}$ 的环都是 $S'_i$ 的环上的\Term{projective-module}{投射模}，\footnote{编者注（53）：一方面，如 VIII，注记 6.8 所述，我们将“自由”改为“\Term{projective-object}{投射}”。另一方面，$S$-\Term{scheme}{概形}\Term{essentially-free}{本质自由}这一概念，与 M. Raynaud 和 L. Gruson 在 [RG71] 中引入和研究的\Term{flat}{平坦}且\Term{pure-scheme}{纯}的 $S$-\Term{scheme}{概形}这一几何概念密切相关，参见后面的增补 6.9。}
则称 $f$ 是\Term{essentially-free}{本质自由}的，也称 $X$ 在 $S$ 上\Term{essentially-free}{本质自由}。
\end{SGAresult}

\begin{SGAresult}{命题}{6.2.2}
a) 若 $X$ 在 $S$ 上\Term{essentially-free}{本质自由}，则它在 $S$ 上\Term{flat}{平坦}；当 $S$ 为\Term{artinian}{Artin}时，逆命题成立。

b) 若 $S$ 是一个域的谱，则任意 $S$-\Term{scheme}{概形}都在 $S$ 上\Term{essentially-free}{本质自由}。

c) 若 $X$ 在 $S$ 上\Term{essentially-free}{本质自由}，则对任意 $S'\to S$，$X'=X\times_S S'$ 在 $S'$ 上\Term{essentially-free}{本质自由}。若 $S'\to S$ \Term{faithfully-flat}{忠实平坦}且\Term{quasi-compact}{拟紧}，则逆命题成立。
\end{SGAresult}
\SourcePage{374}
\begin{proof}
证明直接可得；对于 a) 的逆命题，使用 Artin \Term{local-ring}{局部环}上的\Term{flat-module}{平坦模}是\Term{free-module}{自由模}这一事实。\footnote{编者注（54）：事实上，设 $(A,\mathfrak m)$ 为 Artin \Term{local-ring}{局部环}，$k$ 为其\Term{residue-field}{剩余域}，$M$ 为任意 $A$-模，$(x_i)_{i\in I}$ 为 $M$ 中的一族元素，其像构成 $M/\mathfrak mM$ 的一个 $k$-基。设 $F$ 为以 $(e_i)_{i\in I}$ 为基的\Term{free-module}{自由} $A$-模，$\phi:F\to M$ 为由 $\phi(e_i)=x_i$ 定义的 $A$-\Term{morphism}{态射}。则 $Q=\operatorname{Coker}\phi$ 满足 $Q=\mathfrak mQ$；由于 $\mathfrak m$ \Term{nilpotent}{幂零}，得 $Q=0$。若再假设 $M$ 在 $A$ 上\Term{flat}{平坦}，则 $K=\operatorname{Ker}\phi$ 满足 $K\otimes_A k=0$，即 $K=\mathfrak mK$，因而 $K=0$。}
\end{proof}
引入定义 6.2.1 的理由在于下述\Term{theorem}{定理}。

\begin{SGAresult}{\Term{theorem}{定理}}{6.2.3}
设 $S$ 为一个\Term{scheme}{概形}，$Z$ 为一个\Term{essentially-free}{本质自由}的 $S$-\Term{scheme}{概形}，$Y$ 为 $Z$ 的一个\Term{closed-subscheme}{闭子概形}。考虑由下列条件定义的\Term{functor}{函子}
\[
F=\prod_{Z/S}Y/Z:(\mathrm{Sch})^\circ_{/S}\longrightarrow(\mathrm{Ens}):
\]
当 $Y_{S'}\ne Z_{S'}$ 时 $F(S')=\varnothing$，否则 $F(S')$ 只有一个元素。\footnote{编者注（55）：参见第 II 讲第 1 节，那里将此\Term{functor}{函子}记作 $\prod_{Z/S}Y$；对任意 $S'\to S$，$F(S')=\Gamma(Y_{S'}/Z_{S'})$，在这里，当 $Y_{S'}=Z_{S'}$ 时它等于 $\{\operatorname{id}_{Z_{S'}}\}$，否则为空集。另外，我们增加了断言 (ii)，它用于 6.5.3 的证明。}

(i) 此\Term{functor}{函子}可由 $S$ 的一个\Term{closed-subscheme}{闭子概形} $T$ 表示。

(ii) 若再有 $Y\to Z$ \Term{finitely-presented}{有限表示}，则 $T\to S$ 也\Term{finitely-presented}{有限表示}。
\end{SGAresult}
\begin{proof}
首先注意，所考虑的\Term{functor}{函子}是 (fpqc) \Term{topology}{拓扑}下的一个\Term{sheaf}{层}：因为对每个 $S'$，$F(S')$ 等于 $\varnothing$ 或 $\{\mathrm{pt}\}$，只需验证如下事实：若 $(S_i)$ 是 $S$ 的一个\Term{open-covering}{开覆盖}（相应地，$S'\to S$ 为\Term{faithfully-flat}{忠实平坦}且\Term{quasi-compact}{拟紧}的\Term{morphism}{态射}），且每个 $Y_{S_i}\to Z_{S_i}$（相应地，$Y_{S'}\to Z_{S'}$）都是\Term{isomorphism}{同构}，则 $Y\to Z$ 也是\Term{isomorphism}{同构}。这显然成立（相应地，由 SGA 1，VIII 5.4 或 EGA $\mathrm{IV}_2$，2.7.1 得出）。

此外，由 SGA 1，VIII 1.9，\Term{faithfully-flat}{忠实平坦}且\Term{quasi-compact}{拟紧}的\Term{morphism}{态射}对由\Term{closed-immersion}{闭浸入}\Term{arrow}{箭头}组成的\Term{fibred-category}{纤维范畴}具有\Term{effective-descent}{有效下降}性质。因此，使用 6.2.1 的记号，只需处理 $S=S'_i$ 的情形。

取 $Z$ 的一个\Term{affine-open}{仿射开集}\Term{covering}{覆盖} $(Z_j)$，使 $\mathscr O(Z_j)$ 为 $A=\mathscr O(S)$ 上的\Term{free-module}{自由模}；令 $Y_j=Y\cap Z_j$，并令 $F_j:(\mathrm{Sch})^\circ_{/S}\to(\mathrm{Ens})$ 为由 $(Z_j,Y_j)$ 按照 $F$ 由 $(Z,Y)$ 的定义方式定义的\Term{functor}{函子}。它是\Term{final-functor}{终函子}的一个\Term{subfunctor}{子函子}，且显然 $F=\bigcap_jF_j$。因此，只需证明每个 $F_j$ 可由 $S$ 的一个\Term{closed-subscheme}{闭子概形} $T_j$ 表示（此时 $F$ 便由这些 $T_j$ 的交所成的\Term{closed-subscheme}{闭子概形} $T$ 表示）。故还可假设 $Z$ 为\Term{affine}{仿射}，$Z=\operatorname{Spec}(B)$，其中 $B$ 为一个\Term{projective-module}{投射} $A$-模。设 $L$ 为一个以 $(e_\lambda)_{\lambda\in\Lambda}$ 为基的\Term{free-module}{自由} $A$-模，并且 $B$ 作为 $A$-模是 $L$ 的一个\Term{direct-summand}{直和因子}；令 $\varphi_\lambda:L\to A$ 为关于此基的“坐标”\Term{linear-form}{线性形式}。设 $E$ 为定义 $Z$ 的\Term{closed-subscheme}{闭子概形} $Y$ 的\Term{ideal}{理想} $J\subset B$ 的一个\Term{generator}{生成元}集，$I$ 为 $A$ 中由所有坐标 $\varphi_\lambda(x)$（$x\in E$）生成的\Term{ideal}{理想}。对任意 $A$-\Term{algebra}{代数} $C$，\Term{morphism}{态射} $B\otimes_A C\to(B/J)\otimes_A C$ 为\Term{isomorphism}{同构}，当且仅当对每个 $x\in E$，$x\otimes1$ 在 $L\otimes_A C$ 中的像为零；这又\Term{equivalence}{等价}于 $A\to C$ 的核包含\Term{ideal}{理想} $I$。因此 $T=V(I)=\operatorname{Spec}(A/I)$ 满足所需条件，(i) 得证。

此外，若 $Y\to Z$ \Term{finitely-presented}{有限表示}，则可以取 $E$ 为\Term{finite}{有限}集，于是 $I$ 是 $A$ 的\Term{finite-type}{有限型}\Term{ideal}{理想}，即\Term{closed-immersion}{闭浸入} $T\to S$ \Term{finitely-presented}{有限表示}。
\end{proof}

\SourcePage{375}
\textbf{例 6.2.4。}下面给出一些重要的\Term{functor}{函子}例子，它们可归结为 6.2.3 中的 $\prod_{Z/S}Y/Z$ 型\Term{functor}{函子}；后面需要它们的\Term{representability}{可表性}判据。用 $S$ 表示一个\Term{scheme}{概形}，用 $X,Y,Z$ 等表示 $S$ 上的\Term{scheme}{概形}。

a) 给定一个 $S$-\Term{morphism}{态射}
\[
q:X\longrightarrow\underline{\operatorname{Hom}}_S(Y,Z),\tag{x}
\]
（“$X$ \Term{action}{作用}于 $Y$，取值于 $Z$”），也就是一个\Term{morphism}{态射}
\[
r:X\times_S Y\longrightarrow Z.\tag{xx}
\]
考虑 $Z$ 的一个\Term{subscheme}{子概形} $Z'$，由此有一个\Term{monomorphism}{单态射}
\[
\underline{\operatorname{Hom}}_S(Y,Z')\longrightarrow\underline{\operatorname{Hom}}_S(Y,Z),
\]
使第一个\Term{functor}{函子}成为第二个的\Term{subfunctor}{子函子}。令 $X'$ 为该\Term{subfunctor}{子函子}在\Term{morphism}{态射} $q$ 下的\Term{inverse-image}{逆像}；它是 $X$ 的\Term{subfunctor}{子函子}，其 $X'(T)$ 由满足下述条件的 $x\in X(T)$ 组成：$q(x):Y_T\to Z_T$ 经过 $Z'_T$ \Term{resolution}{分解}。这个\Term{functor}{函子} $X'$ 可如下描述：令 $P=X\times_S Y$，$P'$ 为 $Z'$ 在 $r:P\to Z$ 下的\Term{inverse-image}{逆像}，则有显然的\Term{isomorphism}{同构}
\[
X'\simeq\prod_{P/X}P'/P.\tag{xxx}
\]
因此得到：若 $Y$ 在 $S$ 上\Term{essentially-free}{本质自由}且 $Z'$ 在 $Z$ 中闭，则 $X$ 的\Term{subfunctor}{子函子} $X'$ 可由 $X$ 的一个\Term{closed-subscheme}{闭子概形}表示。\footnote{编者注（56）：下句为新增。}若再有 $Z'\to Z$ \Term{finitely-presented}{有限表示}，则 $X'\to X$ 也\Term{finitely-presented}{有限表示}。

b) 给定 $X$ 以两种方式\Term{action}{作用}于 $Y$，并取值于 $Z$，即给定两个\Term{morphism}{态射}
\[
q_1,q_2:X\rightrightarrows\underline{\operatorname{Hom}}_S(Y,Z),
\]
并令 $X'=\operatorname{Ker}(q_1,q_2)$：这是 $X$ 的\Term{subfunctor}{子函子}，$X'(T)$ 由满足下述条件的 $x\in X(T)$ 组成：两个\Term{morphism}{态射} $q_1(x),q_2(x):Y_T\rightrightarrows Z_T$ 相等。给定 $q_1,q_2$ \Term{equivalence}{等价}于给定一个\Term{morphism}{态射}
\[
q:X\longrightarrow\underline{\operatorname{Hom}}_S(Y,Z\times_S Z),
\]
亦即一个\Term{morphism}{态射} $r:X\times_S Y\to Z\times_S Z$。令 $U=Z\times_S Z$，$U'$ 为 $Z\times_S Z$ 的\Term{diagonal-subscheme}{对角子概形}；则 $X'$ 正是 $\underline{\operatorname{Hom}}_S(Y,U)$ 的\Term{subfunctor}{子函子} $\underline{\operatorname{Hom}}_S(Y,U')$ 在 $q$ 下的\Term{inverse-image}{逆像}。因此它可写成 (xxx) 的\Term{form}{形式}，其中 $P=X\times_S Y$，$P'$ 为\Term{diagonal}{对角线}在 $r$ 下的\Term{inverse-image}{逆像}，即 $r_1,r_2:X\times_S Y\rightrightarrows Z$ 的核。于是可应用 (a)。

因此得到：若 $Y$ 在 $S$ 上\Term{essentially-free}{本质自由}且 $Z$ 在 $S$ 上\Term{separated}{分离}，则 $X$ 的\Term{subfunctor}{子函子} $X'$ 可由 $X$ 的一个\Term{closed-subscheme}{闭子概形}表示。\footnote{编者注（56）：下句为新增。}若再有 $Z\to S$ \Term{locally-finite-type}{局部有限型}，则 $X'\to X$ \Term{finitely-presented}{有限表示}。

\SourcePage{376}
c) 给定一个\Term{morphism}{态射}
\[
q:X\longrightarrow\underline{\operatorname{Hom}}_S(Y,Y),
\]
即“$X$ \Term{action}{作用}于 $Y$”。令 $X'$ 为这个\Term{morphism}{态射}的“核”，即 $X$ 的\Term{subfunctor}{子函子}，其 $X'(T)$ 由使 $q(x):Y_T\to Y_T$ 为恒等的 $x\in X(T)$ 组成。引入第二个\Term{homomorphism}{同态}
\[
q':X\longrightarrow\underline{\operatorname{Hom}}_S(Y,Y),
\]
使“$X$ 平凡地\Term{action}{作用}于 $Y$”，便可对这个\Term{functor}{函子}应用 b)。故有：若 $Y$ 在 $S$ 上\Term{essentially-free}{本质自由}且\Term{separated}{分离}，则 $X$ 的\Term{subfunctor}{子函子} $X'=\operatorname{Ker}q$ 可由 $X$ 的一个\Term{closed-subscheme}{闭子概形}表示。\footnote{编者注（56）：下句为新增。}若再有 $Y\to S$ \Term{locally-finite-type}{局部有限型}，则 $X'\to X$ \Term{finitely-presented}{有限表示}。

d) 在 c) 的条件下，考虑 $Y$ 的“在 $X$ 下\Term{invariant}{不变}”的\Term{subfunctor}{子函子} $Y'$：$Y'(T)$ 由这样的 $y\in Y(T)$ 组成，其对应的\Term{morphism}{态射} $\overline q(y):X_T\to Y_T$ 是“值为 $y$ 的常值 $T$-\Term{morphism}{态射}”。如 c) 那样引入 $q'$，以及对应于 $q,q'$ 的\Term{homomorphism}{同态}
\[
\overline q,\overline q':Y\rightrightarrows\underline{\operatorname{Hom}}_S(X,Y),
\]
可见 $Y'$ 正是 $\operatorname{Ker}(\overline q,\overline q')$，故仍可应用 b)（将 $X,Y$ 的角色对换，并取 $Z=Y$）。

因此，若 $X$ 在 $S$ 上\Term{essentially-free}{本质自由}，$Y$ 在 $S$ 上\Term{separated}{分离}，则 $Y$ 的 $X$-\Term{invariants}{不变元}\Term{subfunctor}{子函子} $Y'$ 可由 $Y$ 的一个\Term{closed-subscheme}{闭子概形}表示。\footnote{编者注（56）：下句为新增。}若再有 $Y\to S$ \Term{locally-finite-type}{局部有限型}，则 $Y'\to Y$ \Term{finitely-presented}{有限表示}。

e) 上述各例中的构造在群论中尤其常见。例如，当 $S$-\Term{group-scheme}{群概形} $G$ \Term{action}{作用}于 $S$-\Term{scheme}{概形} $X$ 时，
\[
q:G\longrightarrow\underline{\operatorname{Aut}}_S(X),
\]
只要 $X$ 在 $S$ 上\Term{essentially-free}{本质自由}且\Term{separated}{分离}，$q$ 的核（“$G$ 中\Term{trivial-action}{平凡作用}的\Term{subgroup}{子群}”）就是 $G$ 的一个\Term{closed-subscheme}{闭子概形}（例 c)）；只要 $G$ 在 $S$ 上\Term{essentially-free}{本质自由}且 $X$ 在 $S$ 上\Term{separated}{分离}，\Term{invariants}{不变元}\Term{subobject}{子对象} $X^G$ 就是 $X$ 的一个\Term{closed-subscheme}{闭子概形}（例 d)）。

设 $Y,Z$ 是 $X$ 的\Term{subscheme}{子概形}。考虑 $G$ 的\Term{subfunctor}{子函子} $\operatorname{Transp}_G(Y,Z)$（“$Y$ 到 $Z$ 的\Term{transporter}{输送子}”），其在 $S$ 上的 $T$ 中取值的点为这样的 $g\in G(T)$：$X_T$ 的相应\Term{automorphism}{自同构}满足 $g(Y_T)\subset Z_T$，即诱导的\Term{morphism}{态射} $Y_T\to X_T$ \Term{decomposable}{可分解}为 $Y_T\to Z_T$。因此，若 $Y$ 在 $S$ 上\Term{essentially-free}{本质自由}且 $Z$ 在 $X$ 中闭，则 $\operatorname{Transp}_G(Y,Z)$ 是 $G$ 的一个\Term{closed-subscheme}{闭子概形}（例 a)）。

也可以考虑 $Y$ 到 $Z$ 的\Term{strict-transporter}{严格输送子}，\footnote{编者注（57）：记作 $\operatorname{Transp}^{\mathrm{str}}_G(Y,Z)$。}其在 $S$ 上的 $T$ 中取值的点为使 $g(Y_T)=Z_T$ 的 $g\in G(T)$；它就是
\[
\operatorname{Transp}_G(Y,Z)\cap\sigma\bigl(\operatorname{Transp}_G(Z,Y)\bigr),
\]
其中 $\sigma$ 为 $G$ 的\Term{inversion-morphism}{求逆态射}。因而，
\SourcePage{377}
若 $Y,Z$ 在 $S$ 上都\Term{essentially-free}{本质自由}且在 $X$ 中闭，则 $Y$ 到 $Z$ 的\Term{strict-transporter}{严格输送子}是 $G$ 的一个\Term{closed-subscheme}{闭子概形}。

一个重要情形是 $X=G$，且 $G$ 以\Term{inner-automorphism}{内自同构}\Term{action}{作用}于自身。若 $H$ 是 $G$ 的一个\Term{subscheme}{子概形}，则 $H$ 到自身的\Term{strict-transporter}{严格输送子}也称为 $H$ 在 $G$ 中的\Term{normalizer}{正规化子}，记作 $\operatorname{Norm}_G H$。所以，若 $H$ 是 $G$ 的一个\Term{closed-group-subscheme}{闭群子概形}且在 $S$ 上\Term{essentially-free}{本质自由}，则 $\operatorname{Norm}_G H$ 可由 $G$ 的一个\Term{closed-group-subscheme}{闭群子概形}表示。

最后，设 $Z$ 为 $G$ 的一个\Term{subscheme}{子概形}。它在 $G$ 中的\Term{centralizer}{中心化子} $\operatorname{Centr}_G(Z)$ 是由 d) 的方法定义的 $G$ 的\Term{group-subfunctor}{群子函子}，这里考虑“$Z$ \Term{action}{作用}于 $G$”，其\Term{action}{作用}由 $G$ 的\Term{action}{作用}诱导。故若 $Z$ 在 $S$ 上\Term{essentially-free}{本质自由}且 $G$ 在 $S$ 上\Term{separated}{分离}，则 $\operatorname{Centr}_G(Z)$ 是 $G$ 的一个\Term{closed-group-subscheme}{闭群子概形}。特别地，若 $G$ 在 $S$ 上\Term{essentially-free}{本质自由}且\Term{separated}{分离}，则 $G$ 的\Term{center}{中心} $C$，即 $\operatorname{Centr}_G(G)$，是 $G$ 的一个\Term{closed-group-subscheme}{闭群子概形}。

当 $S$ 是一个域的谱时，6.2.2 b) 表明，在上面 a) 至 e) 各例中，“\Term{essentially-free}{本质自由}”的条件自动满足，只剩下\Term{separated}{分离}性条件。回顾域上的\Term{group-scheme}{群概形}必定\Term{separated}{分离}（VIA，0.3），便可得到例如以下结果。
\begin{SGAresult}{推论}{6.2.5}
设 $G$ 是域 $k$ 上的一个\Term{group-scheme}{群概形}，$Y,Y'$ 是 $G$ 的两个\Term{subscheme}{子概形}。则：

(i) $Y$ 在 $G$ 中的\Term{centralizer}{中心化子}是 $G$ 的一个\Term{closed-group-subscheme}{闭群子概形}；特别地，$G$ 的\Term{center}{中心} $\operatorname{Centr}_G(G)$ 具有这一性质。

(i$'$) 更一般地，若 $u,v:X\to G$ 是\Term{scheme}{概形}\Term{morphism}{态射}，则 $\operatorname{Transp}_G(u,v)$ 可由 $G$ 的一个\Term{closed-subscheme}{闭子概形}表示。

(ii) 若 $Y$ 是闭的，则\Term{transporter}{输送子} $\operatorname{Transp}_G(Y',Y)$ 是 $G$ 的一个\Term{closed-subscheme}{闭子概形}。若 $Y'$ 也是闭的，则对 $\operatorname{Transp}^{\mathrm{str}}_G(Y',Y)$ 有同样结论。

(iii) 对 $G$ 的任意\Term{group-subscheme}{群子概形}\footnote{编者注（58）：事实上，在域 $k$ 上，$G$ 的任意\Term{group-subscheme}{群子概形}都是闭的，参见 VIA，0.5.2。另外，到 6.2.5 为止，来自 VIII 第 6 节的结果已插入完毕。} $H$，$\operatorname{Norm}_G(H)$ 是 $G$ 的一个\Term{closed-group-subscheme}{闭群子概形}。
\end{SGAresult}

\begin{SGAresult}{推论}{6.2.6}
\footnote{编者注（59）：此处插入本推论（参见第 XII 讲，6.1），后面将用到它。另外，从 6.3 起回到 VIB 的原文。}
设 $k$ 为一个域，$G$ 为一个\Term{connected}{连通}的 $k$-\Term{algebraic-group}{代数群}。则 $\operatorname{Centr}_G(G)$ 可由 $G$ 的一个\Term{closed-group-subscheme}{闭群子概形} $Z$ 表示，且 $G/Z$ 为一个\Term{affine}{仿射} $k$-\Term{algebraic-group}{代数群}。
\end{SGAresult}
\begin{proof}
第一条断言当然已包含在 6.2.5 (i) 中，不过我们将看到，它也可从第二条断言的证明中得出。事实上，$G$ 通过\Term{adjoint-representation}{伴随表示}\Term{action}{作用}于\Term{finite-dimensional}{有限维} $k$-\Term{vector-space}{向量空间} $V_n=\mathscr O_{G,e}/\mathfrak m_e^{n+1}$（其中 $\mathfrak m_e$ 为 $\mathscr O_{G,e}$ 的\Term{maximal-ideal}{极大理想}）；记 $K_n$ 为\Term{morphism}{态射} $\rho_n:G\to\operatorname{GL}(V_n)$ 的核。由 VIA 5.4.1，$\rho_n$ 诱导\Term{closed-immersion}{闭浸入} $G/K_n\hookrightarrow\operatorname{GL}(V_n)$，故每个 $G/K_n$ 都是\Term{affine}{仿射}的。由于 $G$ 为\Term{noetherian}{Noether}，$K_n$ 的交 $K$ 等于其中一个 $K_n$，因而 $G/K$ 为\Term{affine}{仿射}。
\SourcePage{378}
另一方面，记 $Z$ 为 $G$ 的\Term{center}{中心}，显然 $Z\subset K$。证明 $Z=K$。记 $\widehat{\mathscr O}_{G,e}$ 为 $\mathscr O_{G,e}$ 关于 $\mathfrak m_e$-进\Term{topology}{拓扑}的\Term{completion}{完备化}，$\widehat S$ 为它的谱，并令 $S=\operatorname{Spec}\mathscr O_{G,e}$。由于 $\widehat S\to S$ \Term{faithfully-flat}{忠实平坦}，而两个\Term{morphism}{态射}
\[
K\times_k S\longrightarrow S,\qquad
(g,x)\longmapsto gxg^{-1}\quad\text{及}\quad(g,x)\longmapsto x
\]
在\Term{base-change}{基变换} $\widehat S\to S$ 后相等，所以它们相等，即 $K$ 平凡地\Term{action}{作用}于 $\mathscr O_{G,e}$。由 6.2.4 e)，$G$ 在 $K$ 下的\Term{invariants}{不变元}\Term{subobject}{子对象} $G^K$（它就是 $\operatorname{Centr}_G(K)$）是 $G$ 的一个\Term{closed-subscheme}{闭子概形}，因而由 $\mathscr O_G$ 的一个\Term{quasi-coherent}{拟凝聚}\Term{ideal}{理想} $\mathscr I$ 定义。由于 $G^K$ 包含 $S=\operatorname{Spec}\mathscr O_{G,e}$，且 $\mathscr I$ 为\Term{finite-type}{有限型}（因为 $G$ 为\Term{noetherian}{Noether}），存在 $e$ 的一个\Term{open-neighborhood}{开邻域} $U$，使 $\mathscr I|_U=0$。于是\Term{subgroup}{子群} $G^K=\operatorname{Centr}_G(K)$ 包含 $U$，故也包含 $U\cdot U$；由于 $G$ \Term{irreducible}{不可约}，后者等于 $G$（VIA 0.5）。因此 $\operatorname{Centr}_G(K)=G$，故 $K\subset Z$，从而 $Z=K$。
\end{proof}

\textbf{注记 6.3。}设 $k$ 为一个\Term{algebraically-closed}{代数闭}域，$G$ 为一个 $k$-\Term{group}{群}，$H$ 为 $G$ 的一个\Term{group-subscheme}{群子概形}；假设 $G,H$ 都在 $k$ 上\Term{finite-type}{有限型}且\Term{reduced}{约化}。则 $\operatorname{Norm}_G H$（相应地，$\operatorname{Centr}_G H$）可由 $G$ 的一个\Term{group-subscheme}{群子概形}表示，其\Term{adjunction}{伴随}的\Term{reduced-subscheme}{约化子概形}正是 Bible 意义下 $H$ 在 $G$ 中的\Term{normalizer}{正规化子}（相应地，\Term{centralizer}{中心化子}）。

\begin{SGAresult}{命题}{6.4}
设 $G$ 为一个 $S$-\Term{group}{群}，$u:X\to G$ 为一个 $S$-\Term{scheme}{概形}\Term{monomorphism}{单态射}。令 $T=\operatorname{Transp}_G(X,X)$。下列条件\Term{equivalence}{等价}：

(i) $T$ 是 $G$ 的一个\Term{group-subfunctor}{群子函子}。

(ii) $T=\operatorname{Transp}^{\mathrm{str}}_G(X,X)=\operatorname{Norm}_G X$。

在以下两个情形的每一个中，这些条件均成立：

a) $X$ 在 $S$ 上\Term{finitely-presented}{有限表示}。

b) $T$ 可由 $S$ 上的一个\Term{finitely-presented}{有限表示}\Term{scheme}{概形}表示。
\end{SGAresult}
\begin{proof}
条件 (i) 与 (ii) \Term{equivalence}{等价}，因为对任意\Term{morphism}{态射} $S'\to S$ 及任意 $t,t'\in T(S')$，有 $tt'\in T(S')$，并且 $\operatorname{Transp}^{\mathrm{str}}_G(X,X)=T\cap c(T)$（参见 6.1 (ii)）。

先考虑 a)。若 $t\in T(S)$，则 $\operatorname{int}(t)$ 为 $X$ 到自身的一个\Term{monomorphism}{单态射}，故为 $X$ 的一个 $S$-\Term{automorphism}{自同构}（EGA $\mathrm{IV}_4$，17.9.6），所以 $t$ 属于 $\operatorname{Transp}^{\mathrm{str}}_G(X,X)$，a) 得证。

在 b) 中，显然 $\mu(T\times_S T)\subset T$，断言便由下述\Term{lemma}{引理}得出。
\end{proof}

\begin{SGAresult}{\Term{lemma}{引理}}{6.4.2}
\footnote{编者注（60）：这里保留原文的编号：没有 6.4.1。}
设 $G$ 为一个\Term{finitely-presented}{有限表示}的 $S$-\Term{scheme}{概形}，带有一个\Term{associative-law}{结合律}（EGA $0_{\mathrm{III}}$，8.2.5 的意义下）。假设对每个 $S$-\Term{scheme}{概形} $S'$ 和每个 $g\in G(S')$，\Term{set}{集合} $G(S')$ 中由 $g$ 给出的左\Term{right-translation}{右平移}都是\Term{injectivity}{单射}，且 $G(S)\ne\varnothing$。则 $G$ 是一个 $S$-\Term{group}{群}。
\end{SGAresult}
\SourcePage{379}
\begin{proof}
只需证明，对任意 $S$-\Term{scheme}{概形} $S'$，\Term{set}{集合} $G(S')$ 都是一个群。由假设立即可知，对每个 $g\in G(S')$，$G_{S'}$ 中由 $g$ 给出的左\Term{right-translation}{右平移}都是 $G_{S'}$ 到自身的 $S$-\Term{monomorphism}{单态射}。由于 $G$ 在 $S$ 上\Term{finitely-presented}{有限表示}，这些都是 $S$-\Term{automorphism}{自同构}（EGA $\mathrm{IV}_4$，17.9.6）。因此，\Term{set}{集合} $G(S')$ 中由 $g$ 给出的左\Term{right-translation}{右平移}都是\Term{bijection}{双射}，问题归结为下述\Term{lemma}{引理}。
\end{proof}
\begin{SGAresult}{\Term{lemma}{引理}}{6.4.3}
设 $G$ 为一个\Term{nonempty}{非空}\Term{set}{集合}，带有满足左\Term{right-translation}{右平移}都是\Term{bijection}{双射}的\Term{associative-law}{结合律}。则 $G$ 是一个群。
\end{SGAresult}
证明留给读者。

\begin{SGAresult}{定义}{6.5}
设 $G$ 为一个 $S$-\Term{group}{群}，$H$ 为一个 $S$-\Term{functor}{函子}，$u:H\to G$ 为一个\Term{monomorphism}{单态射}。

(i) 称\Term{functor}{函子} $\operatorname{Centr}_G H$ 的\Term{neutral-component}{单位连通分支}为 $H$ 在 $G$ 中的\Term{connected-centralizer}{连通中心化子}，记作 $\operatorname{Centr}^0_G H$（参见 3.1 与 6.1 (iii)）。

(i$'$) 对任意\Term{morphism}{态射} $u:X\to G$，类似定义\Term{functor}{函子} $\operatorname{Centr}^0(u)$（参见 6.2 (iv)）。

(ii) 当每个 $s\in S$ 处的 $u_s$ 都是\Term{closed-immersion}{闭浸入}时，称\Term{functor}{函子} $\operatorname{Norm}_G H$ 的\Term{neutral-component}{单位连通分支}为 $H$ 在 $G$ 中的\Term{connected-normalizer}{连通正规化子}，记作 $\operatorname{Norm}^0_G H$。
\end{SGAresult}
\textbf{注意。}由 1.4.2，当 $H$ 是一个 $S$-\Term{group-scheme}{群概形}，$G,H$ 在 $S$ 上\Term{locally-finite-type}{局部有限型}，且 $u$ 是一个\Term{quasi-compact}{拟紧}的 $S$-\Term{group}{群}\Term{morphism}{态射}时，(ii) 的假设成立。

\begin{SGAresult}{命题}{6.5.1}
设 $G$ 为一个在 $S$ 上\Term{locally-finitely-presented}{局部有限表示}且\Term{quasi-separated}{拟分离}的 $S$-\Term{group}{群}，$H$ 为一个\Term{smooth}{光滑}且\Term{connected}{纤维连通}的 $S$-\Term{group}{群}，$u:H\to G$ 为一个\Term{monomorphism}{单态射}。设 $N$ 为 $H$ 在 $G$ 中的\Term{normalizer}{正规化子}（参见 6.1）。由下面的 6.5.5，$N$ 可由 $G$ 的一个\Term{closed-group-subscheme}{闭群子概形}表示，且在 $G$ 上\Term{finitely-presented}{有限表示}。在此条件下，以下条件\Term{equivalence}{等价}：

(i) \Term{canonical}{典范}\Term{morphism}{态射} $H\to N$ 是\Term{open-immersion}{开浸入}。

(ii) $N^0=H$（参见 3.10）。

(iii) 对每个 $s\in S$，$H_s=(N_s)^0$。
\end{SGAresult}
\begin{proof}
由于 $H^0=H$，由\Term{lemma}{引理} 3.10.1，条件 (i) 蕴含 (ii)。另一方面，显然 (ii) 蕴含 (iii)，因为 $H_s=(N^0)_s=(N_s)^0$。

最后证明 (iii) 蕴含 (i)。由于对每个 $s\in S$ 都有 $H_s=(N^0)_s$，故 $H_s\to N_s$ 是\Term{open-immersion}{开浸入}。又因为 $H,N$ 在 $S$ 上\Term{locally-finitely-presented}{局部有限表示}，且 $H$ 在 $S$ 上\Term{flat}{平坦}，由 EGA $\mathrm{IV}_4$，17.9.5，$H\to N$ 是\Term{open-immersion}{开浸入}。
\end{proof}
为方便读者，下面收入第 XI 讲的结果 6.8 至 6.11。\footnote{编者注（61）：此处插入的 6.5.2 至 6.5.5 取自第 XI 讲，6.8 至 6.11。事实上，A. Grothendieck 已在 XI 第 6 节开头的注中建议：“本节不使用（XI 的）第 3、4、5 节的结果；它的自然位置应在 VIB 中。”}

\SourcePage{380}
\begin{SGAresult}{\Term{theorem}{定理}}{6.5.2}
设 $X$ 为 $S$ 上的一个\Term{smooth}{光滑}\Term{scheme}{概形}，其\Term{geometrically-irreducible}{纤维几何不可约}。\footnote{编者注（62）：一方面，根据证明，我们把原文的“\Term{connected}{连通}”改为“\Term{irreducible}{不可约}”。另一方面，由 [RG71] I，3.3.4 (iii) 及 4.1.1，只须假设 $X$ 在 $S$ 上\Term{flat}{平坦}，其\Term{geometrically-irreducible}{纤维几何不可约}且没有\Term{embedded-component}{嵌入分支}。}

(i) 对 $X$ 的任意\Term{closed-subscheme}{闭子概形} $Y$，\Term{functor}{函子} $\prod_{X/S}Y/X$ 可由 $S$ 的一个\Term{closed-subscheme}{闭子概形} $T$ 表示。

(ii) 若再有 $Y\to X$ \Term{finitely-presented}{有限表示}，则 $T\to S$ 也\Term{finitely-presented}{有限表示}。
\end{SGAresult}
\begin{proof}
由于 $f:X\to S$ \Term{faithfully-flat}{忠实平坦}且\Term{locally-finitely-presented}{局部有限表示}，它对 (fpqc) \Term{topology}{拓扑}是\Term{covering-morphism}{覆盖态射}。另一方面，$T=\prod_{X/S}Y/X$ 显然是 (fpqc) \Term{topology}{拓扑}下 $S$ 的一个\Term{subsheaf}{子层}，所以 $T$ 能否由 $S$ 的\Term{closed-subscheme}{闭子概形}表示这一问题，在 $S$ 上对于 (fpqc) \Term{topology}{拓扑}是\Term{local-scheme}{局部}的；\footnote{编者注（63）：例如参见第 VIII 讲的增补 1.7。}判断 $T$ 是否在 $S$ 上\Term{finitely-presented}{有限表示}也是如此（参见 EGA $\mathrm{IV}_2$，2.7.1）。因此，经\Term{base-change}{基变换} $S'\to S$，其中 $S'=X$，可归结到 $X$ 在 $S$ 上具有一个\Term{section}{截面} $\varepsilon$ 的情形。还可假设 $S$ 为\Term{affine}{仿射}，因而\Term{quasi-compact}{拟紧}。此时有如下结果。
\end{proof}

\begin{SGAresult}{推论}{6.5.3}
在 6.5.2 的条件下，假设 $S$ \Term{quasi-compact}{拟紧}，$X\to S$ 有一个\Term{section}{截面} $\varepsilon$，且 $Y\to X$ \Term{finitely-presented}{有限表示}。则存在整数 $n\geq0$，使
\[
\prod_{X/S}Y/X=\prod_{X_n/S}Y_n/X_n,
\]
其中 $X_n$ 为\Term{immersion}{浸入} $\varepsilon:S\to X$ 的第 $n$ 阶\Term{infinitesimal-neighborhood}{无穷小邻域}，且 $Y_n=Y\cap X_n$。
\end{SGAresult}
当 $Y$ 在 $X$ 上\Term{finitely-presented}{有限表示}时，这个推论确实蕴含 6.5.2：由于 $X$ 在 $S$ 上\Term{smooth}{光滑}，$X_n$ 在 $S$ 上\Term{finite-morphism}{有限}且\Term{locally-free}{局部自由}，因而更是“\Term{essentially-free}{本质自由}”的（参见 6.2.1），故由 6.2.3，$\prod_{X_n/S}Y_n/X_n$ 可由 $S$ 的一个\Term{closed-subscheme}{闭子概形} $T$ 表示，且在 $S$ 上\Term{finitely-presented}{有限表示}。

\begin{proof}
先在 $S$ 为\Term{noetherian}{Noether}的情形证明 6.5.3（从而证明 6.5.2）。设 $T_n=\prod_{X_n/S}Y_n/X_n$，则 $T_n$ 构成 $S$ 的\Term{closed-subscheme}{闭子概形}的一个递减序列；由于 $S$ 为\Term{noetherian}{Noether}，该序列最终\Term{stable}{稳定}。记
\[
R=\bigcap_{n\geq0}\prod_{X_n/S}Y_n/X_n
\]
为 $n$ \Term{large-relative-to-cardinal}{足够大}时它们的共同值。显然 $T\subset R$，只需证明 $R\subset T$。经\Term{base-change}{基变换} $R\to S$，可归结到 $R=S$，即每个 $n$ 都有 $Y_n=X_n$，也就是每个 $n$ 都有 $Y\supset X_n$，此时要证明 $T=S$，即 $Y=X$。

由 $Y\supset X_n$ 对所有 $n$ 成立，并利用 $X$ \Term{locally-noetherian}{局部 Noether}，可知在 $\varepsilon(S)$ 每一点的\Term{neighborhood}{邻域}内，$Y$ 都是 $X$ 的一个诱导\Term{open-subscheme}{开子概形}。\footnote{编者注（64）：例如参见 6.2.6 的证明。}所以存在 $X$ 的一个诱导\Term{open-subset}{开集} $U$，包含 $\varepsilon(S)$，且 $U\subset Y$。由 EGA $\mathrm{IV}_3$，11.10.10，由于 $X/S$ 的\Term{fiber}{纤维}\Term{integral}{整}，$U$ 在 $X$ 中\Term{schematically-dense}{概形稠密}；于是（因 $Y$ 是包含 $U$ 的\Term{closed-subscheme}{闭子概形}）$Y=X$。这就在 $S$ 为\Term{noetherian}{Noether}时证明了 6.5.3，从而证明了 6.5.2。

\SourcePage{381}
一般情形通过归结到上述情形处理。对每个 $s\in S$，存在 $s$ 的一个\Term{affine-open}{仿射开}\Term{neighborhood}{邻域} $U$，及 $\varepsilon(s)$ 的一个\Term{affine-open}{仿射开}\Term{neighborhood}{邻域} $V$，使 $f(V)\subset U$。于是 $f(V)$ 是包含在 $U$ 中的 $s$ 的一个\Term{open-neighborhood}{开邻域}。若 $S'$ 是包含在 $\varepsilon^{-1}(V)\cap f(V)$ 中的 $s$ 的一个\Term{affine-open}{仿射开}\Term{neighborhood}{邻域}，并令 $X'=V\cap f^{-1}(S')$，则 $X',S'$ 分别是 $X,S$ 的\Term{affine-open}{仿射开集}，且 $X'/S'$ 有一个\Term{section}{截面}。由于 6.5.2 与 6.5.3 具有\Term{local-character}{局部性}，可假设 $S=S'$。于是 $X'$ 为包含 $\varepsilon(S)$ 的 $X$ 的一个\Term{affine-open}{仿射开集}。由于每个\Term{fiber}{纤维} $X_s$ 都假设为\Term{irreducible}{不可约}，$X'_s$ 在 $X_s$ 中\Term{schematically-dense}{概形稠密}；故由 EGA $\mathrm{IV}_3$，11.10.10，$X'$ 在 $X$ 中\Term{schematically-dense}{概形稠密}。同样，对任意\Term{base-change}{基变换} $S_1\to S$，$X'\times_S S_1$ 在 $X\times_S S_1$ 中\Term{schematically-dense}{概形稠密}。

由此得
\[
\prod_{X/S}Y/X=\prod_{X'/S}Y'/X',\qquad Y'=Y\cap X'.
\]
这使我们归结到 $X=X'$ 的情形，所以可以假设 $S,X$ 都为\Term{affine}{仿射}。此外，若 $X=\operatorname{Spec}(B)$，$J$ 为定义 $Y$ 的 $B$ 的\Term{ideal}{理想}，则 $J$ 是其\Term{finite-type}{有限型}\Term{subideal}{子理想}的\Term{inductive-limit}{归纳极限}，故 $Y$ 是 $X$ 的一些在 $X$ 上\Term{finitely-presented}{有限表示}的\Term{closed-subscheme}{闭子概形} $Y_i$ 的交。因此
\[
\prod_{X/S}Y/X=\bigcap_i\prod_{X/S}Y_i/X,
\]
所以，为证明 6.5.2，只须处理 $Y$ 在 $X$ 上\Term{finitely-presented}{有限表示}且 $S,X$ 都为\Term{affine}{仿射}的情形。

此时 $S$ 上的 $X,Y$ 是由某个\Term{noetherian}{Noether}\Term{scheme}{概形} $S_0$ 上的类似情形 $X_0,Y_0$ 经\Term{base-change}{基变换} $S\to S_0$ 得到的，因而归结到已处理的 $S$ 为\Term{noetherian}{Noether}的情形。这就完成了 6.5.2 和 6.5.3 的证明。
\end{proof}

\begin{SGAresult}{推论}{6.5.4}
设 $X$ 为一个\Term{smooth}{光滑}、\Term{finitely-presented}{有限表示}且\Term{connected}{纤维连通}的 $S$-\Term{group-scheme}{群概形}，$Y$ 为一个在 $S$ 上\Term{finitely-presented}{有限表示}的\Term{group-scheme}{群概形}，$i:Y\to X$ 为一个 $S$-\Term{group-scheme}{群概形}\Term{monomorphism}{单态射}。

(i) 则 $\prod_{X/S}Y/X$ 可由 $S$ 的一个\Term{finitely-presented}{有限表示}\Term{closed-subscheme}{闭子概形}表示。

(ii) 若再假设 $S$ \Term{quasi-compact}{拟紧}，则对充分大的 $n$，有
\[
\prod_{X/S}Y/X=\prod_{X_n/S}Y_n/X_n,
\]
其中 $X_n$ 表示\Term{unit-section}{单位截面} $\varepsilon:S\to X$ 的第 $n$ 阶\Term{infinitesimal-neighborhood}{无穷小邻域}，且 $Y_n=X_n\cap Y$。
\end{SGAresult}
\begin{proof}
证明实质上与 6.5.3 的证明相同。\footnote{编者注（65）：下面详细展开了原文的论证，并增加了 EGA IV 的参考。}一方面，$i$ \Term{locally-finitely-presented}{局部有限表示}（参见 EGA $\mathrm{IV}_1$，1.4.3）。另一方面，$Y,X$ 的\Term{unit-section}{单位截面}诱导\Term{bijection}{双射}\Term{immersion}{浸入} $S\to Y_n$ 与 $S\to X_n$，因而诱导 $S_{\mathrm{r\text{é}d}}$ 与 $(Y_n)_{\mathrm{r\text{é}d}}$、$(X_n)_{\mathrm{r\text{é}d}}$ 的\Term{isomorphism}{同构}。所以 $i_n$ \Term{quasi-compact}{拟紧}，从而\Term{finite-type}{有限型}，并有如下\Term{commutative-diagram}{交换图}：
\[
\begin{tikzcd}
(Y_n)_{\mathrm{r\text{é}d}}\arrow[r,"\tau"]\arrow[dr,"\sigma"']&Y_n\arrow[d,"i_n"]\\
&X_n.
\end{tikzcd}
\]
\SourcePage{382}
其中 $\sigma,\tau$ 是\Term{closed-immersion}{闭浸入}，且 $\tau$ 为\Term{surjectivity}{满射}。由于 $i_n$ \Term{separated}{分离}（它是\Term{monomorphism}{单态射}），由此 $i_n$ \Term{proper}{固有}（参见 EGA II，5.4.3）。所以 $i_n$ 为\Term{finitely-presented}{有限表示}的\Term{proper}{固有}\Term{monomorphism}{单态射}，因而为\Term{closed-immersion}{闭浸入}（参见 EGA $\mathrm{IV}_3$，8.11.5）。于是，由前面已用过的 6.2.3，$\prod_{X_n/S}Y_n/X_n$ 可由 $S$ 的一个在 $S$ 上\Term{finitely-presented}{有限表示}的\Term{closed-subscheme}{闭子概形} $T_n$ 表示。因此，只须在另假设 $S$ 为\Term{affine}{仿射}的情形下证明 6.5.4 的最后一条断言。

仍可立即归结到 $S$ 为\Term{noetherian}{Noether}的情形，此时要证明 $R=T$（采用 6.5.3 证明中的记号），亦即 $Y\supset X_n$ 对所有 $n$ 成立蕴含 $Y=X$。这个假设说明，$i:Y\to X$ 在 $Y$ 的\Term{unit-section}{单位截面}各点处\Term{etale}{平展}，因而 $Y$ 在\Term{unit-section}{单位截面}各点处在 $S$ 上\Term{smooth}{光滑}。于是，$Y$ 中在 $S$ 上\Term{smooth}{光滑}的点组成的\Term{open-subset}{开集} $U$ 是 $Y$ 的一个诱导\Term{open-subgroup}{开子群}（参见 2.3）。由 2.5，$\tau:U\to X$ 是\Term{etale}{平展}\Term{monomorphism}{单态射}，从而是\Term{open-immersion}{开浸入}。但 $X$ 的\Term{fiber}{纤维}\Term{connected}{连通}，而\Term{algebraic-group}{代数群}的每个\Term{open-subgroup}{开子群}也是闭的，因此 $\tau$ 是满的\Term{open-immersion}{开浸入}，即\Term{isomorphism}{同构}。故 $U=X$，更有 $Y=X$，6.5.4 得证。
\end{proof}
如 6.2.4 e) 那样，由 6.5.4 可\Term{pushout}{推出}：

\begin{SGAresult}{推论}{6.5.5}
设 $G$ 为一个在 $S$ 上\Term{locally-finite-type}{局部有限型}且\Term{quasi-separated}{拟分离}的 $S$-\Term{group-scheme}{群概形}，$H$ 为一个在 $S$ 上\Term{smooth}{光滑}、\Term{finitely-presented}{有限表示}且\Term{connected}{纤维连通}的\Term{group-scheme}{群概形}，$i:H\to G$ 为一个 $S$-\Term{group}{群}\Term{monomorphism}{单态射}。则：

a) $\operatorname{Centr}_G(H)$ 与 $\operatorname{Norm}_G(H)$ 可由 $G$ 的\Term{closed-subscheme}{闭子概形}表示，且在 $G$ 上\Term{finitely-presented}{有限表示}。

a$'$) 同样，对任意\Term{finitely-presented}{有限表示}的 $S$-\Term{group-scheme}{群概形}\Term{monomorphism}{单态射} $j:K\to G$，$\operatorname{Transp}_G(H,K)$ 可由 $G$ 的一个\Term{closed-subscheme}{闭子概形}表示，且在 $G$ 上\Term{finitely-presented}{有限表示}。

b) 若 $G$ \Term{quasi-compact}{拟紧}，则存在整数 $n\geq0$，使得（以 $H_n$ 表示 $H$ 的\Term{unit-section}{单位截面}的第 $n$ 阶\Term{infinitesimal-neighborhood}{无穷小邻域}）
\[
\begin{gathered}
\operatorname{Centr}_G(H)=\operatorname{Centr}_G(H_n),\qquad
\operatorname{Norm}_G(H)=\operatorname{Norm}_G(H_n),\\
\operatorname{Transp}_G(H,K)=\operatorname{Transp}_G(H_n,K)
=\operatorname{Transp}_G(H_n,K_n).
\end{gathered}
\]
\end{SGAresult}
\begin{proof}
\footnote{编者注（66）：下面的论证作了详细展开。另外，到这里来自 XI，6.8 至 6.11 的插入结束，即下面从 6.6 起回到 VIB 的原文。}
首先注意，对 $G$ 的假设说明\Term{monomorphism}{单态射} $H\to G$ \Term{finitely-presented}{有限表示}（EGA $\mathrm{IV}_1$，1.2.4 与 1.4.3），\Term{diagonal-immersion}{对角浸入} $\Delta_{G/S}:G\to G\times_S G$ 也一样（同上，1.4.3.1）。所以，取 $H=K$，就将 $\operatorname{Norm}_G(H)$ 的情形归结到\Term{transporter}{输送子}的情形。利用 6.2.4 e)，我们将对以 $G$ 为基\Term{scheme}{概形}的\Term{group-scheme}{群概形} $X=H_G=H\times_S G$ 及以下\Term{group-subscheme}{群子概形} $Y$ 应用 6.5.4。

对 $\operatorname{Transp}_G(H,K)$，取 $Y$ 为 $K_G$ 在 $G$-\Term{group}{群}\Term{morphism}{态射}
\[
H\times_S G\longrightarrow G\times_S G,\qquad(h,g)\longmapsto(ghg^{-1},g)
\]
下的\Term{inverse-image}{逆像}。对
\SourcePage{383}
$\operatorname{Centr}_G(H)$，取 $Y$ 为 $(G\times_S G)_G$ 的\Term{diagonal-subgroup}{对角子群} $\Delta_{G/S}(G)_G$ 在 $G$-\Term{group}{群}\Term{morphism}{态射}
\[
H\times_S G\longrightarrow G\times_S G\times_S G,\qquad
(h,g)\longmapsto(h,ghg^{-1},g)
\]
下的\Term{inverse-image}{逆像}。
\end{proof}

\begin{SGAresult}{定义}{6.6}
设 $G$ 为一个 $S$-\Term{group-functor}{群函子}，$H$ 为一个 $S$-\Term{group-subfunctor}{群子函子}。若 $\operatorname{Norm}_G H=G$（相应地，$\operatorname{Centr}_G H=G$；相应地，对任意 $S$-\Term{scheme}{概形} $T$ 和任意\Term{automorphism}{自同构} $a\in\operatorname{Aut}_{T\text{-gr.}}(G_T)$，有 $a(H_T)\subset H_T$），则称 $H$ 在 $G$ 中\Term{normal-subgroup}{正规}（相应地，\Term{central-subgroup}{中心}；相应地，\Term{characteristic-subgroup}{特征}）。换言之，对任意 $S$-\Term{scheme}{概形} $T$，$G(T)$ 的\Term{subgroup}{子群} $H(T)$ 在 $G(T)$ 中\Term{normal-subgroup}{正规}（相应地，在 $G(T)$ 中\Term{central-subgroup}{中心}；相应地，在 $G_T$ 的任意\Term{automorphism}{自同构}下\Term{invariant}{不变}）。
\end{SGAresult}
\textbf{注意。}若 $H$ 为\Term{central-subgroup}{中心}（相应地，\Term{characteristic-subgroup}{特征}）\Term{subgroup}{子群}，则它\Term{normal-subgroup}{正规}。

\textbf{注记 6.7。}设 $G,H$ 为两个 $S$-\Term{group}{群}，$u:H\to G$ 为一个\Term{monomorphism}{单态射}。$H$ 在 $G$ 中\Term{normal-subgroup}{正规}（相应地，\Term{central-subgroup}{中心}）的充要条件是\Term{morphism}{态射}
\[
\mu\circ\bigl(c\circ\operatorname{pr}_2,\mu\circ(u\times\operatorname{id}_G)\bigr):H\times_S G\longrightarrow G
\]
（对任意 $g\in G(S')$ 与 $h\in H(S')$ 由 $(h,g)\mapsto g^{-1}hg$ 定义）经过 $u$ \Term{resolution}{分解}（相应地，等于 $u\circ\operatorname{pr}_1$）。$H$ 在 $G$ 中为\Term{characteristic-subgroup}{特征}\Term{subgroup}{子群}的充要条件是：对任意 $S$-\Term{scheme}{概形} $T$ 及 $G_T$ 的任意 $T$-\Term{group}{群}\Term{automorphism}{自同构} $a$，$a\circ u_{(T)}$ 经过 $u_{(T)}$ \Term{resolution}{分解}。

\textbf{例 6.8。}设 $G$ 为一个 $S$-\Term{group-functor}{群函子}。则 $\operatorname{Centr}G$ 为\Term{characteristic-subgroup}{特征}且\Term{central-subgroup}{中心}\Term{subgroup}{子群}。若对每个 $s\in S$，$G_s$ 都\Term{representable}{可表}，则 $G^0$ 为\Term{characteristic-subgroup}{特征}\Term{subgroup}{子群}。这由定义和 3.3 得出。

\textbf{6.9。}\footnote{编者注（67）：本小节为新增。}
在 [RG71]，I 3.3.3 中，作者引入了\Term{pure-scheme}{纯} $S$-\Term{scheme}{概形}这一几何概念，它对于 $S$ 上的\Term{etale}{平展}\Term{topology}{拓扑}是\Term{local-scheme}{局部}的；精确定义参见该处。这里只指出以下几点：

a)（同上，\Term{theorem}{定理} 3.3.5）若 $B$ 是一个\Term{flat}{平坦}、\Term{finitely-presented}{有限表示}的 $A$-\Term{algebra}{代数}，则 $B$ 为\Term{projective-module}{投射} $A$-模，当且仅当 $\operatorname{Spec}(B)$ 在 $\operatorname{Spec}(A)$ 上\Term{pure-scheme}{纯}。

b) 因而，若 $X\to S$ \Term{locally-finitely-presented}{局部有限表示}、\Term{flat}{平坦}且\Term{pure-scheme}{纯}，则 $X$ 在 $S$ 上\Term{essentially-free}{本质自由}。

c)（同上，3.3.4 (iii)）若 $X\to S$ \Term{locally-finite-type}{局部有限型}、\Term{flat}{平坦}，其\Term{geometrically-irreducible}{纤维几何不可约}且没有\Term{embedded-component}{嵌入分支}，则 $X$ 在 $S$ 上\Term{pure-scheme}{纯}。

由于域上的每个\Term{locally-finite-type}{局部有限型}\Term{group-scheme}{群概形}都是 Cohen--Macaulay 的（参见 VIA，1.1.1），因而没有\Term{embedded-component}{嵌入分支}，特别得到：

d) 每个在 $S$ 上\Term{locally-finitely-presented}{局部有限表示}、\Term{flat}{平坦}且\Term{connected}{纤维连通}的 $S$-\Term{group-scheme}{群概形} $G$，都在 $S$ 上\Term{pure-scheme}{纯}，从而在 $S$ 上\Term{essentially-free}{本质自由}。

利用上述 (b) 和 (d) 的结果，可以重新考察 6.2.4 e) 中的全部命题。

\SourcePage{384}
\SGAsection{7}{生成子群；交换子群}
本节中，$k$ 表示一个固定的域。

\begin{SGAresult}{命题}{7.1}
设 $G$ 为一个 $k$-\Term{group}{群}，$(X_i)_{i\in I}$ 为一族\Term{geometrically-reduced}{几何约化}的 $k$-\Term{scheme}{概形}；\footnote{编者注（68）：这里及下文将不常用的术语“\Term{separable}{可分}”改为通常的术语“\Term{geometrically-reduced}{几何约化}”，参见 EGA $\mathrm{IV}_2$，4.6.2。}对每个 $i\in I$，给定一个 $k$-\Term{morphism}{态射} $f_i:X_i\to G$。

(i) 存在包含每个 $f_i$ 的最小的 $G$ 的\Term{closed-group-subscheme}{闭} $k$-\Term{group-subscheme}{群子概形}，记作 $\Gamma_G((f_i)_{i\in I})$。它是一个\Term{geometrically-reduced}{几何约化}的 $k$-\Term{scheme}{概形}，因而当 $G$ 在 $k$ 上\Term{locally-finite-type}{局部有限型}时，它\Term{smooth}{光滑}（1.3.1）。

(ii) 令 $X=\coprod_{i\in I}X_i$，$f:X\to G$ 为在每个 $X_i$ 上的\Term{restriction}{限制}为 $f_i$ 的\Term{morphism}{态射}。令 $X^1=X\coprod X$，$f^1:X^1\to G$ 为在两份 $X$ 上的\Term{restriction}{限制}分别为 $f$ 和 $c\circ f$ 的\Term{morphism}{态射}。对每个 $n>1$，令
\[
X^n=X^1\times_k X^{n-1},\qquad
f^n=\mu\circ(f^1\times_k f^{n-1}):X^n\longrightarrow G.
\]
则 $\Gamma_G((f_i)_{i\in I})$ 是 $G$ 中以 $\bigcup_{n\geq1}f^n(X^n)$ 的\Term{closure}{闭包}为底\Term{space}{空间}的\Term{reduced-subscheme}{约化子概形}。

(iii) 对任意 $k$-\Term{scheme}{概形} $S$，$\Gamma_G((f_i)_{i\in I})_S$ 是 $G_S$ 的包含每个 $f_{i,S}:X_S\to G_S$ 的最小\Term{closed-group-subscheme}{闭群子概形}。

(iii$'$) 而且，$\Gamma_G((f_i)_{i\in I})_S$ 是 $G_S$ 的包含每个 $f_{i,S}$ 的最小\Term{group-subscheme}{群子概形}。\footnote{编者注（69）：原文陈述 (iii$'$) 时另假设 $G$ 在 $k$ 上\Term{locally-finite-type}{局部有限型}；由 VIA，0.5.2，这一假设可以去掉。}
\end{SGAresult}
\begin{proof}
首先注意，按 (ii) 定义 $X,f$，则为证明 (i)、(iii) 和 (iii$'$)，可归结到 $I$ 只有一个元素的情形。

设 $H$ 为 $G$ 中以 $\overline{\bigcup_{n\geq1}f^n(X^n)}$ 为\Term{underlying-set}{底集}的\Term{reduced-subscheme}{约化子概形}。则\Term{morphism}{态射}族 $f^n:X^n\to H$ \Term{schematically-dominant}{概形支配}（参见 EGA $\mathrm{IV}_3$，11.10.4），所以 $G$ 的每个包含所有 $f^n$ 的\Term{closed-subscheme}{闭子概形}也包含 $H$。此外，由同上 11.10.7，$H$ \Term{geometrically-reduced}{几何约化}。故为证明 (i) 和 (ii)，只须证明 $H$ 是 $G$ 的一个\Term{group-subscheme}{群子概形}，即 $c$ 在 $H$ 上的\Term{restriction}{限制}和 $\mu$ 在 $H\times_k H$ 上的\Term{restriction}{限制}都经过\Term{injectivity}{单射} $H\to G$ \Term{resolution}{分解}。

由于 $H$ \Term{geometrically-reduced}{几何约化}，$H\times_k H$ \Term{reduced}{约化}，只须在\Term{set}{集合}意义下验证 $c(H)\subset H$ 及 $\mu(H\times_k H)\subset H$。由 EGA $\mathrm{IV}_3$，11.10.6，所有 $f^n_{(H)}(X^n\times_k H)$ 的并在 $H\times_k H$ 中\Term{schematically-dense}{概形稠密}。同样，对每个 $n\geq1$，所有 $f^m_{(X^n)}(X^n\times_k X^m)$（$m\geq1$）的并在 $X^n\times_k H$ 中\Term{schematically-dense}{概形稠密}。故只须证明
\[
\mu\bigl(f^n_{(H)}(f^m_{(X^n)}(X^n\times_k X^m))\bigr)\subset H,
\qquad c(f^n(X^n))\subset H.
\]
而
\[
\begin{aligned}
\mu\bigl(f^n_{(H)}(f^m_{(X^n)}(X^n\times_k X^m))\bigr)
&=\mu\bigl((f^n\times_k f^m)(X^n\times_k X^m)\bigr)\\
&=f^{n+m}(X^{n+m})\subset H;
\end{aligned}
\]
并且，由于 $c(f^1(X^1))\subset f^1(X^1)$，对任意 $n$ 有 $c(f^n(X^n))\subset f^n(X^n)\subset H$。这就证明了 (i) 和 (ii)。

\SourcePage{385}
现在证明 (iii)。设 $G'$ 为 $G_S$ 中包含 $f_S$ 的一个\Term{closed-group-subscheme}{闭群子概形}；要证明 $G'$ 包含 $H_S$，\Term{equivalence}{等价}地，$H_S=H_S\times_{G_S}G'$。令 $H'_S=H_S\times_{G_S}G'=G'\cap H_S$。由于 $G',H_S$ 都包含所有 $f^n_S$，$H'_S$ 也如此。由 EGA $\mathrm{IV}_3$，11.10.6，由于 $f^n:X^n\to H$ 这一族\Term{schematically-dominant}{概形支配}，$f^n_S:X^n_S\to H_S$ 这一族也\Term{schematically-dominant}{概形支配}；因此，由同上 11.10.1 c)，包含每个 $f^n_S$ 的 $H'_S$ 等于 $H_S$。(iii) 得证。

最后证明，$H_S$ 是 $G_S$ 中包含 $f_S$ 的最小\Term{group-subscheme}{群子概形}（未必预先假设为闭）。

设 $G'$ 为 $G_S$ 中包含 $f_S$ 的一个\Term{group-subscheme}{群子概形}。仍需证明，令 $H'_S=H_S\times_{G_S}G'$，则 $H'_S=H_S$。只须证明 $H'_S$ 在 $H_S$ 中闭，并应用 (iii)。为此只须证明 $H_S,H'_S$ 有相同的\Term{underlying-set}{底集}；更强地，只须证明对每个 $s\in S$，$H_s$ 等于
\[
H'_s:=H'_S\times_S\kappa(s)=H_s\times_{G_s}G'_s.
\]
由 VIA，0.5.2，$G_s$ 的 $\kappa(s)$-\Term{group-subscheme}{群子概形} $G'_s$ 在 $G_s$ 中闭。故 $H'_s$ 在 $H_s$ 中闭；将前述论证用于 $H'_s,H_s$ 及 $f^n_s$，便知 $H'_s=H_s$。
\end{proof}

\begin{SGAresult}{推论}{7.1.1}
\footnote{编者注（70）：增加本推论是为指出 7.1 的这一特殊情形。}
设 $G$ 为一个\Term{locally-finite-type}{局部有限型} $k$-\Term{group}{群}，$A,B$ 为 $G$ 的两个\Term{smooth}{光滑}、\Term{finite-type}{有限型} $k$-\Term{subgroup}{子群}，$i_A$（相应地，$i_B$）为 $A$（相应地，$B$）到 $G$ 的包含\Term{morphism}{态射}。假设 $B$ \Term{normalize}{正规化} $A$，则 $A\cdot B=\Gamma_G(i_A,i_B)$。
\end{SGAresult}
\begin{proof}
事实上，设 $H=A\rtimes B$ 为 $A,B$ 的\Term{semidirect-product}{半直积}（参见 I，2.3.5），它是一个\Term{smooth}{光滑}、\Term{finite-type}{有限型} $k$-\Term{group}{群}。群\Term{morphism}{态射} $u:H\to G$，$(a,b)\mapsto ab$，\Term{quasi-compact}{拟紧}，故由 1.2，$u(H)=A\cdot B$ 是 $G$ 的一个\Term{reduced}{约化}\Term{closed-subscheme}{闭子概形}，并且是\Term{category}{范畴} $(\mathrm{Sch}/k)_{\mathrm{r\text{é}d}}$ 中的一个群。由 EGA $\mathrm{IV}_3$，11.10.7，$A\cdot B=u(H)$ \Term{geometrically-reduced}{几何约化}（因为 $H$ 如此），所以它是 $G$ 的一个\Term{closed-subgroup}{闭子群}。由于显然 $A\cdot B\subset\Gamma_G(i_A,i_B)$，推论随之成立。
\end{proof}

\textbf{定义与注记 7.2。}\footnote{编者注（71）：把原文的 (viii) 移到 (i$'$)，并把 (vi)、(vii) 单列为下面的推论 7.2.1 和定义 7.2.2。}
(i) 给定一个 $k$-\Term{group}{群} $G$、一族\Term{geometrically-reduced}{几何约化}的 $k$-\Term{scheme}{概形} $(X_i)_{i\in I}$，以及对每个 $i\in I$ 给定的 $k$-\Term{morphism}{态射} $f_i:X_i\to G$，称包含每个 $f_i$ 的最小\Term{closed-group-subscheme}{闭群子概形}为由族 $(f_i)_{i\in I}$ \Term{generated-subgroup}{生成的 $G$ 的闭群子概形}，在本节记作 $\Gamma_G((f_i)_{i\in I})$。若 $X$ 为 $G$ 的一个在 $k$ 上\Term{geometrically-reduced}{几何约化}的\Term{subscheme}{子概形}，$f$ 为\Term{immersion}{浸入} $X\hookrightarrow G$，则将 $\Gamma_G(f)$ 写作 $\Gamma_G(X)$。

(i$'$) 采用 7.1 (ii) 的记号，有时记 $\Gamma'_G(f)=\bigcup_{n\geq1}f^n(X^n)$。注意，$\Gamma'_G(f)$ 是 $G$ 中对\Term{group-law}{群运算}\Term{stable}{稳定}的一个\Term{subset}{子集}（3.0 的意义下）。

(ii) 显然，若 $X_1,X_2$ 是两个\Term{geometrically-reduced}{几何约化}的 $k$-\Term{scheme}{概形}，$f_1:X_1\to G$ 与 $f_2:X_2\to G$ 为两个 $k$-\Term{morphism}{态射}，且\Term{set}{集合} $\overline{f_1(X_1)},\overline{f_2(X_2)}$ 相等，则 $\Gamma_G(f_1)=\Gamma_G(f_2)$。

\SourcePage{386}
(iii) 设 $E$ 为 $G$ 的一个\Term{subset}{子集}，使得 $G$ 的\Term{reduced-subscheme}{约化子概形} $\overline E$ \Term{geometrically-reduced}{几何约化}。称\Term{group-subscheme}{群子概形} $\Gamma_G(i)$（其中 $i$ 为\Term{reduced-subscheme}{约化子概形} $\overline E$ 到 $G$ 的\Term{injectivity}{单射}）为 $E$ \Term{generated-subgroup}{生成的 $G$ 的闭群子概形}，记作 $\Gamma_G(E)$。

(iv) 由 VIA，0.5.2，$G$ 的每个\Term{group-subscheme}{群子概形}都是闭的，\footnote{编者注（72）：此处去掉了 $G$ 在 $k$ 上\Term{locally-finite-type}{局部有限型}的假设。}所以将“生成的\Term{closed-group-subscheme}{闭群子概形}”简称为“\Term{generated-subgroup}{生成的群子概形}”。

(v) 设 $X$ 为一个\Term{geometrically-reduced}{几何约化}的 $k$-\Term{scheme}{概形}，$f:X\to G$ 为一个 $k$-\Term{morphism}{态射}。假设 $f(X)$ 包含 $G$ 的\Term{identity-element}{单位元} $e$。令 $X'{}^1=X\times_k X$，$f'{}^1=\mu\circ(f\times_k(c\circ f))$，并对 $n>1$ 令
\[
X'{}^n=X'{}^1\times_k X'{}^{n-1},\qquad
f'{}^n=\mu\circ(f'{}^1\times_k f'{}^{n-1}).
\]
则 $\Gamma_G(f)$ 是 $G$ 中以 $\bigcup_{n\geq1}f'{}^n(X'{}^n)$ 的\Term{closure}{闭包}为底\Term{space}{空间}的\Term{reduced-subscheme}{约化子概形}。

事实上，回顾 7.1 (ii) 的记号：$X^1=X\amalg X$，$X^n=X^1\times_k X^{n-1}$（$n\geq2$），有以下包含关系，其中第一个来自假设 $e\in f(X)$：
\[
f^n(X^n)\subset f'{}^n(X'{}^n)\subset f^{2n}(X^{2n}),\qquad\forall n\geq1.
\]
这还表明，存在整数 $n$ 使 $f^n(X^n)=\Gamma_G(f)$ 的充要条件是，存在整数 $m$ 使 $f'{}^m(X'{}^m)=\Gamma_G(f)$。

由注记 7.2 (v) 得到：
\begin{SGAresult}{推论}{7.2.1}
设 $X$ 为一个\Term{geometrically-reduced}{几何约化}且\Term{geometrically-connected}{几何连通}的 $k$-\Term{scheme}{概形}，$f:X\to G$ 为一个 $k$-\Term{morphism}{态射}，使 $f(X)$ 包含 $G$ 的\Term{identity-element}{单位元}。则 $k$-\Term{group}{群} $\Gamma_G(f)$ \Term{connected}{连通}，从而\Term{irreducible}{不可约}。
\end{SGAresult}
\begin{proof}
此时每个 $X'{}^n$ 都\Term{connected}{连通}，故所有 $f'{}^n(X'{}^n)$（它们都包含\Term{identity-element}{单位元}）的并\Term{connected}{连通}，其\Term{closure}{闭包} $\Gamma_G(f)$ 也\Term{connected}{连通}。于是，由 VIA，2.4（当 $G$，因而 $\Gamma_G(f)$，在 $k$ 上\Term{locally-finite-type}{局部有限型}时）及 2.6.5 (iii)（一般情形），$\Gamma_G(f)$ \Term{irreducible}{不可约}。
\end{proof}

\begin{SGAresult}{定义}{7.2.2}
设 $G$ 为一个 $k$-\Term{group}{群}。

a) 设 $A,B$ 为 $G$ 的两个\Term{geometrically-reduced}{几何约化}的 $k$-\Term{group-subscheme}{群子概形}。称由\Term{morphism}{态射} $\nu:A\times_k B\to G$ 生成的 $G$ 的\Term{closed-group-subscheme}{闭群子概形}为 $A,B$ 在 $G$ 中的\Term{commutator-subgroup}{交换子群概形}，记作 $(A,B)_G$，或简记为 $(A,B)$；这里 $\nu$ 对任意 $k$-\Term{scheme}{概形} $S$ 及 $a\in A(S),b\in B(S)$，由 $(a,b)\mapsto aba^{-1}b^{-1}$ 定义。

b) 假设 $G$ 在 $k$ 上\Term{geometrically-reduced}{几何约化}。称群 $(G,G)$ 为 $G$ 的\Term{derived-group}{导群}，记作 $\mathrm D(G)$。\footnote{编者注（73）：将原文的记号 $\operatorname{Der}(G)$ 改掉，以免与\Term{derivation}{导子}\Term{space}{空间}混淆。}
\end{SGAresult}
\textbf{注意。}$G$ \Term{commutative-group}{交换}的充要条件是 $\mathrm D(G)$ 为单位 $k$-\Term{group}{群}。

\SourcePage{387}
\textbf{回顾 7.3.0。}\footnote{编者注（74）：增加了这些回顾，并相应修改及详细展开了 7.3 的表述。}
回顾，若 $\phi:Y\to Z$ 为一个 $S$-\Term{scheme}{概形}\Term{morphism}{态射}，则\Term{image-presheaf}{像预层} $\phi(Y)$ 为如下 $S$-\Term{functor}{函子}：它将 $S$ 上的任意 $S'$ 对应到 $Z(S')$ 的\Term{subset}{子集} $\phi(Y(S'))$。注意，若 $T$ 为 $Z$ 的一个\Term{subscheme}{子概形}，且\Term{presheaf}{预层}的包含 $\phi(Y)\hookrightarrow Z$ 经过 $T$ \Term{resolution}{分解}，即对任意 $S$-\Term{morphism}{态射} $h:S'\to Y$ 有 $\phi\circ h\in T(S')$，则在\Term{set}{集合}意义下 $\phi(Y)\subset T$（取 $h=\operatorname{id}_Y$）；当 $Y$ \Term{reduced}{约化}时逆命题成立，因为此时 $\phi$ 经过 $T$ \Term{resolution}{分解}。

还应回顾，由 6.2，若 $H$ 为 $G$ 的一个\Term{closed-group-subscheme}{闭} $k$-\Term{group-subscheme}{群子概形}，则 $\operatorname{Centr}_G H$ 与 $\operatorname{Norm}_G H$ 可由 $G$ 的\Term{closed-group-subscheme}{闭} $k$-\Term{group-subscheme}{群子概形}表示。

\begin{SGAresult}{推论}{7.3}
设 $G$ 为一个 $k$-\Term{group}{群}，$X$ 为一个\Term{geometrically-reduced}{几何约化}的 $k$-\Term{scheme}{概形}，$f:X\to G$ 为一个 $k$-\Term{morphism}{态射}。

(i) 设 $S$ 为一个 $k$-\Term{scheme}{概形}，$u$ 为 $S$-\Term{group}{群} $G_S$ 的一个\Term{endomorphism}{自同态}。

(a) 若在\Term{set}{集合}意义下 $u(f_S(X_S))\subset\Gamma_G(f)_S$，则\Term{morphism}{态射} $u:\Gamma_G(f)_S\to G_S$ 经过 $\Gamma_G(f)_S$ \Term{resolution}{分解}。

(b) 若 $u$ 为 $S$-\Term{group}{群} $G_S$ 的一个\Term{automorphism}{自同构}，且在\Term{set}{集合}意义下 $u(f_S(X_S))=f_S(X_S)$，则 $u$ 诱导 $\Gamma_G(f)_S$ 的一个\Term{automorphism}{自同构}。特别地，若元素 $g\in G(S)$ 在\Term{set}{集合}意义下满足 $\operatorname{int}(g)(f_S(X_S))=f_S(X_S)$，则 $g\in\operatorname{Norm}_{G_S}(\Gamma_G(f)_S)(S)$。

(c) 若 $u\circ f_S=f_S$，则 $u$ 在 $G_S$ 的\Term{subgroup}{子群} $\Gamma_G(f)_S$ 上的\Term{restriction}{限制}为恒等。特别地，若元素 $g\in G(S)$ 满足 $\operatorname{int}(g)(f_S)=f_S$，则 $g\in\operatorname{Centr}_{G_S}(\Gamma_G(f)_S)(S)$。

(ii) 设 $H$ 为 $G$ 的一个\Term{group-subscheme}{群子概形}，则 $H$ \Term{centralize}{中心化}（相应地，\Term{normalize}{正规化}）$\Gamma_G(f)$，当且仅当对任意 $S\to\operatorname{Spec}k$ 及 $h\in H(S)$ 有 $\operatorname{int}(h)\circ f_S=f_S$（相应地，$\operatorname{int}(h)(f_S(X_S))\subset\Gamma_G(f)_S$）；也就是，对任意 $S'\to S$ 及 $x\in X(S')$，$G(S')$ 的元素 $h_{S'},f(x)$ 交换（相应地，$h_{S'}f(x)h_{S'}^{-1}\in\Gamma_G(f)(S')$）。

(iii) 特别地，再设 $Y$ 为一个\Term{geometrically-reduced}{几何约化}的 $k$-\Term{scheme}{概形}，$\phi:Y\to G$ 为一个 $k$-\Term{morphism}{态射}。假设对任意 $k$-\Term{scheme}{概形} $S'$ 及任意 $x\in X(S'),y\in Y(S')$，$G(S')$ 的元素 $\phi(y),f(x)$ 交换（相应地，$\operatorname{int}(\phi(y))(f(x))\in\Gamma_G(f)(S')$）。则 $\Gamma_G(\phi)$ 是 $\operatorname{Centr}_G\Gamma_G(f)$（相应地，$\operatorname{Norm}_G\Gamma_G(f)$）的一个 $k$-\Term{subgroup}{子群}。

(iv) 若 $g$ 为 $G$ 的一个 $k$-点，则 $k$-\Term{group}{群} $\Gamma_G(\{g\})$ \Term{commutative-group}{交换}。

(v) 设 $A,B$ 为 $G$ 的两个在 $k$ 上\Term{geometrically-reduced}{几何约化}的\Term{group-subscheme}{群子概形}。若 $A,B$ 都\Term{normal-subgroup}{正规}（相应地，都是\Term{characteristic-subgroup}{特征}\Term{subgroup}{子群}），则 $(A,B)$ 也如此。
\end{SGAresult}
\begin{proof}
(i) 先证 (a)。令 $\Gamma_S=\Gamma_G(f)_S$。则 $u^{-1}(\Gamma_S)$ 是 $G_S$ 的一个\Term{closed-group-subscheme}{闭} $S$-\Term{group-subscheme}{群子概形}，故由 7.1 (iii)，只需证明 $f_S$ 经过 $u^{-1}(\Gamma_S)$ \Term{resolution}{分解}。由于 $u(f_S(X_S))\subset\Gamma_S$ 且 $X_S$ \Term{reduced}{约化}，$u\circ f_S$ 经过 $\Gamma_S$ \Term{resolution}{分解}，故 $f_S$ 经过 $u^{-1}(\Gamma_S)$ \Term{resolution}{分解}。(a) 得证。

\SourcePage{388}
于是，将 (a) 用于 $u$ 和 $u^{-1}$，便得到 (b) 的第一条断言（事实上只须假设 $u(f_S(X_S))\subset\Gamma_S$ 及 $u^{-1}(f_S(X_S))\subset\Gamma_S$）；第二条断言是取 $u=\operatorname{int}(g)$ 的特殊情形。

证明 (c)。考虑 $S$-\Term{group}{群}\Term{morphism}{态射} $G_S\to G_S\times_S G_S$，$g\mapsto(u(g),g)$，并以 $H$ 表示\Term{diagonal}{对角线}的\Term{inverse-image}{逆像}，它是 $G_S$ 的一个\Term{group-subscheme}{群子概形}。由于 $G$ 在 $k$ 上\Term{separated}{分离}（VIA 0.3），$G_S$ 在 $S$ 上\Term{separated}{分离}，故 $H$ 在 $G_S$ 中闭。由假设，$H$ 包含 $f_S$，故由 7.1 (iii) 它包含 $\Gamma_G(f)_S$。这证明了 (c) 的第一条断言；第二条是取 $u=\operatorname{int}(g)$ 的特殊情形。

证明 (ii)。令 $\Gamma=\Gamma_G(f)$，记 $i$ 为包含\Term{morphism}{态射} $\Gamma\hookrightarrow G$。$H$ 包含在 $N=\operatorname{Norm}_G(\Gamma)$ 中，当且仅当对每个 $k$-\Term{scheme}{概形} $S$ 和 $h\in H(S)$，有 $\operatorname{int}(h)(\Gamma_S)\subset\Gamma_S$（将这个条件用于 $h,h^{-1}$，即得 $\operatorname{int}(h)(\Gamma_S)=\Gamma_S$）；由 (i)(a)，其充要条件正是 $\operatorname{int}(h)(f_S(X_S))\subset\Gamma_S$。

同样，$H$ 包含在 $C=\operatorname{Centr}_G(\Gamma)$ 中，当且仅当对每个 $k$-\Term{scheme}{概形} $S$ 及 $h\in H(S)$，有 $\operatorname{int}(h)\circ i_S=i_S$；由 (i)(c)，这又当且仅当 $\operatorname{int}(h)(f_S)=f_S$。(ii) 得证。

证明 (iii)。由 (ii)，假设说明 $\phi$ 经过 $C$（相应地，$N$）\Term{resolution}{分解}；由 6.2，后者是 $G$ 的一个\Term{closed-subgroup}{闭子群}，故由 7.1 (i)，它包含 $\Gamma_G(f)$。

在 (iii) 中取 $f,\phi$ 均为 $g$ 所定义的\Term{closed-immersion}{闭浸入} $\operatorname{Spec}k\hookrightarrow G$，即得 (iv)：此时对每个 $k$-\Term{scheme}{概形} $S$，$X(S),Y(S)$ 都只有一个点，且该点由 $f$（相应地，$\phi$）映到 $g$。

最后证明 (v)。设 $\nu$ 为\Term{morphism}{态射} $G\times G\to G$，$(g,h)\mapsto ghg^{-1}h^{-1}$，$\nu'$ 为它在 $A\times B$ 上的\Term{restriction}{限制}；由定义（7.2.2），$(A,B)=\Gamma_G(\nu')$。设 $S$ 为一个 $k$-\Term{scheme}{概形}，$u$ 为 $G_S$ 的一个\Term{inner-automorphism}{内自同构}（相应地，一个 $S$-\Term{group}{群}\Term{automorphism}{自同构}）。有 $u\circ\nu_S=\nu_S\circ(u\times u)$。另一方面，假设说明 $u$ 诱导 $A_S$ 的一个\Term{automorphism}{自同构} $u_1$（相应地，$B_S$ 的一个\Term{automorphism}{自同构} $u_2$），故
\[
u\bigl(\nu'_S(A_S\times_S B_S)\bigr)
=\nu'_S\bigl(u_1(A_S)\times_S u_2(B_S)\bigr)
=\nu'_S(A_S\times_S B_S).
\]
于是由 (i)(b)，$u$ 诱导 $(A,B)_S$ 的一个\Term{automorphism}{自同构}。(v) 得证。
\end{proof}

\begin{SGAresult}{命题}{7.4}
设 $G$ 为一个\Term{locally-finite-type}{局部有限型} $k$-\Term{group}{群}，$X$ 为一个\Term{finite-type}{有限型}、\Term{geometrically-reduced}{几何约化}且\Term{geometrically-connected}{几何连通}的 $k$-\Term{scheme}{概形}，$f:X\to G$ 为一个 $k$-\Term{morphism}{态射}，使 $f(X)$ 包含 $G$ 的\Term{identity-element}{单位元} $e$。则采用 7.1 (ii) 的记号，存在整数 $N$，使在\Term{set}{集合}意义下
\[
f^N(X^N)=\Gamma_G(f).
\]
\end{SGAresult}
\begin{proof}
由 7.1 (iii) 及 EGA $\mathrm{IV}_2$，2.6.1，可假设 $k$ \Term{algebraically-closed}{代数闭}。由推论 7.2.1，可假设 $G=G^0$；最后，采用 7.2 (v) 的记号，只须证明存在整数 $N$，使 $f'{}^N(X'{}^N)=\Gamma_G(f)$。

\textbf{第一种情形。}假设 $X$ \Term{irreducible}{不可约}。则 $\overline{f'{}^n(X'{}^n)}$ 构成\Term{space}{空间} $G$ 中\Term{irreducible}{不可约}闭集的一个递增序列；由于 $G=G^0$ 在 $k$ 上\Term{finite-type}{有限型}（VIA 2.4），$G$ 为\Term{noetherian}{Noether}。
\SourcePage{389}
所以这个序列最终\Term{stable}{稳定}，存在整数 $m$，使 $\overline{f'{}^m(X'{}^m)}=\Gamma_G(f)$。

此外，由于 $X,G$ 在 $k$ 上\Term{finite-type}{有限型}，\Term{morphism}{态射} $f'{}^n$ 在 $k$ 上\Term{finite-type}{有限型}。因此 $f'{}^m(X'{}^m)$ 在 $G$ 中\Term{constructible}{可构造}（EGA $\mathrm{IV}_1$，1.8.5），故包含其\Term{closure}{闭包} $\Gamma_G(f)$ 的一个\Term{open-subset}{开集} $U$（EGA $0_{\mathrm{IV}}$，9.2.3）。由 VIA 0.5，有
\[
\Gamma_G(f)\subset U\cdot U\subset f'{}^{2m}(X'{}^{2m})\subset\Gamma_G(f),
\]
故 $f'{}^{2m}(X'{}^{2m})=\Gamma_G(f)$。

\textbf{第二种情形。}假设 $X$ 恰有两个\Term{irreducible-component}{不可约分支} $A_1,A_2$。由于 $X$ \Term{connected}{连通}且 $k$ \Term{algebraically-closed}{代数闭}，存在 $a\in(A_1\cap A_2)(k)$。于是四个\Term{irreducible}{不可约}\Term{subset}{子集} $A_i\times_k A_j$（$i,j=1,2$）\Term{covering}{覆盖} $X\times_k X$，且每个\Term{subset}{子集}在\Term{morphism}{态射} $f'{}^1=\mu\circ(f\times_k(c\circ f))$ 下的像都包含 $e$。以 $f'{}^1_{ij}$ 表示 $f'{}^1$ 在\Term{reduced-subscheme}{约化子概形} $A_i\times_k A_j$ 上的\Term{restriction}{限制}，令
\[
\begin{gathered}
Y=(A_1\times_k A_1)\times_k(A_1\times_k A_2)
\times_k(A_2\times_k A_2)\times_k(A_2\times_k A_1),\\
g=\mu\circ\Bigl(\bigl(\mu\circ(f'{}^1_{11}\times_k f'{}^1_{12})\bigr)
\times_k\bigl(\mu\circ(f'{}^1_{22}\times_k f'{}^1_{21})\bigr)\Bigr).
\end{gathered}
\]
则 $Y$ \Term{irreducible}{不可约}、\Term{reduced}{约化}且\Term{finite-type}{有限型}，所以由刚才的证明，存在整数 $m$，使 $g'{}^m(Y'{}^m)=\Gamma_G(g)$。对任意 $n\geq1$，有
\[
f'{}^n(X'{}^n)\subset g'{}^n(Y^n)\subset f'{}^{4n}(X'{}^{4n}),
\]
故 $\Gamma_G(f)=\Gamma_G(g)$ 且 $f'{}^{4m}(X'{}^{4m})=\Gamma_G(f)$。

\textbf{一般情形。}对 $X$ 的\Term{irreducible-component}{不可约分支}数作归纳（由于 $X$ 在 $k$ 上\Term{finite-type}{有限型}，从而为\Term{noetherian}{Noether}，这个数\Term{finite}{有限}）。假设命题在 $X$ 至多有 $r-1$ 个\Term{irreducible-component}{不可约分支}时成立，并假设现有 $r$ 个，即 $A_1,\ldots,A_r$。$e$ 属于某个 $A_i$ 的像；不妨设 $e\in f(A_1)$。令 $Y=\Gamma_G(f_r)$，其中 $f_r$ 为 $f$ 在\Term{reduced-subscheme}{约化子概形} $X_r=A_1\cup\cdots\cup A_{r-1}$ 上的\Term{restriction}{限制}（假设 $A_i$ 的编号选得使此\Term{scheme}{概形}\Term{connected}{连通}，这总能做到）。则 $Y$ 是 $G$ 的一个闭的、\Term{reduced}{约化}的\Term{subgroup}{子群}，并由推论 7.2.1 知它\Term{irreducible}{不可约}。

令 $Z=\overline{f(A_r)}$，$T=Y\cup Z$，赋予\Term{reduced}{约化}\Term{closed-subscheme}{闭子概形}结构，并令 $i$ 为 $T$ 到 $G$ 的\Term{injectivity}{单射}。显然（7.2 (ii)）$\Gamma_G(i)=\Gamma_G(f)$，且 $T$ \Term{connected}{连通}（因为 $X$ \Term{connected}{连通}，$A_1\cup\cdots\cup A_{r-1}$ 与 $A_r$ 有公共点 $a$，于是 $Y,Z$ 有公共点 $f(a)$）。此外，$e\in T$，且由于 $Y,Z$ 都\Term{irreducible}{不可约}，$T$ 至多有两个\Term{irreducible-component}{不可约分支}。由归纳假设，存在整数 $m$，使 $f_r'{}^m(X_r'{}^m)=\Gamma_G(f_r)=Y$。由于 $X=X_r\cup A_r$，且 $Z=\overline{f(A_r)}$ 包含在 $\overline{f'{}^m(X'{}^m)}$ 中（因为 $e\in f(A_1)$），故
\[
f(X)\subset Y\cup Z\subset\overline{f'{}^m(X'{}^m)}.
\]
另一方面，由于 $T$ 至多有两个\Term{irreducible-component}{不可约分支}，前面已证明存在整数 $p$，使 $i'{}^p(T'{}^p)=\Gamma_G(i)=\Gamma_G(f)$。又 $T\subset\overline{f'{}^m(X'{}^m)}$，故 $\overline{f'{}^{mp}(X'{}^{mp})}=\Gamma_G(f)$；如第一种情形所示，这个等式蕴含 $f'{}^{2mp}(X'{}^{2mp})=\Gamma_G(f)$。
\end{proof}

\SourcePage{390}
\begin{SGAresult}{\Term{lemma}{引理}}{7.5}
设 $S$ 为一个\Term{scheme}{概形}，$G$ 为一个 $S$-\Term{group-scheme}{群概形}，$X$ 为一个 $S$-\Term{prescheme}{预概形}，$f:X\to G$ 为一个 $S$-\Term{morphism}{态射}。假设 $X,G$ 在 $S$ 上\Term{locally-finitely-presented}{局部有限表示}，并且对每个 $s\in S$ 及 $G_s$ 的每个\Term{maximal-point}{极大点} $g$，存在 $X$ 的一个点 $x$，使 $f(x)=g$ 且 $f$ 在 $x$ 处\Term{flat}{平坦}。则\Term{morphism}{态射} $\mu\circ(f\times_S f):X\times_S X\to G$ 对 (fppf) \Term{topology}{拓扑}是\Term{covering-morphism}{覆盖态射}。
\end{SGAresult}
\begin{proof}
\footnote{编者注（75）：下面详细展开了原文的论证。}
由 EGA $\mathrm{IV}_3$，11.3.1，$X$ 中使 $f$ \Term{flat}{平坦}的点组成一个\Term{open-subset}{开集} $V$，且 $f|_V$ 是\Term{open-morphism}{开态射}，故 $U=f(V)$ 是 $G$ 的一个\Term{open-subset}{开集}。又由假设，对每个 $s\in S$，$U\cap G_s$ 在 $G_s$ 中\Term{dense}{稠密}。记 $\phi$ 为 $\mu\circ(f\times_S f)$ 在 $V\times_S V$ 上的\Term{restriction}{限制}。

只须证明 $\phi$ 对 (fppf) \Term{topology}{拓扑}是\Term{covering-morphism}{覆盖态射}；为此，只须证明 $\phi$ \Term{faithfully-flat}{忠实平坦}且\Term{finitely-presented}{有限表示}（参见 IV，6.3.1 (iv)）。$\phi$ 等于复合
\[
V\times_S V\xrightarrow{\,f|_V\times f|_V\,}U\times_S U\xrightarrow{\,\mu\,}G,
\]
其中第一个\Term{morphism}{态射}\Term{faithfully-flat}{忠实平坦}且\Term{locally-finitely-presented}{局部有限表示}，因为 $f|_V$ 如此。故只须证明 $\mu$ 在 $U\times_S U$ 上的\Term{restriction}{限制}也如此。

由于 $G\to S$ \Term{locally-finitely-presented}{局部有限表示}且\Term{flat}{平坦}，$\mu:G\times_S G\to G$ 也如此（它\Term{isomorphism}{同构}于由 $G\to S$ 经\Term{base-change}{基变换} $G\to S$ 得到的\Term{morphism}{态射}，参见 VIA 0.1），因而诱导的\Term{morphism}{态射} $U\times_S U\to G$ 也如此。为证明后者\Term{surjectivity}{满射}，只须逐\Term{fiber}{纤维}考察；由于对每个 $s\in S$，$U\cap G_s$ 是 $G_s$ 的\Term{dense-open-set}{稠密开集}，\Term{surjectivity}{满射}性由 VIA 0.5 得出。
\end{proof}

\textbf{注记 7.6.0。}\footnote{编者注（76）：此定义为新增；在原文中，它包含在命题 7.6 的陈述里。}
设 $S$ 为一个\Term{scheme}{概形}，$H$ 为一个 $S$-\Term{group}{群}，$f:X\to H$ 为一个 $S$-\Term{scheme}{概形}\Term{morphism}{态射}。$H$ 中由 $f$ 的像生成的\Term{group-subpresheaf}{群子预层}，记作 $\langle\operatorname{Im}f\rangle$，是 $H$ 的如下 $k$-\Term{group-subfunctor}{群子函子}：它将每个 $S$-\Term{scheme}{概形} $S'$ 对应到 $H(S')$ 中由 $f(X(S'))$ \Term{generated-subgroup}{生成的子群}。由于该\Term{subgroup}{子群}的每个元素都可写成\Term{finite}{有限}乘积 $f(x_1)^{\varepsilon_1}\cdots f(x_n)^{\varepsilon_n}$，其中 $n\in\mathbf N^*$、$x_i\in X(S')$、$\varepsilon_i=\pm1$，所以若记 $X^1=X\coprod X$，并如 7.1 定义\Term{morphism}{态射} $f^n:X^n\to H$，则 $\langle\operatorname{Im}f\rangle$ 正是\Term{morphism}{态射}
\[
f^\infty:X^\infty\longrightarrow H
\]
的\Term{image-presheaf}{像预层}，其中 $X^\infty=\coprod_{n\geq1}X^n$，且 $f^\infty:X^\infty\to H$ 为在每个 $X^n$ 上\Term{restriction}{限制}为 $f^n$ 的 $S$-\Term{morphism}{态射}。

\begin{SGAresult}{命题}{7.6}
\footnote{编者注（77）：原文在下述特殊情形的假设下陈述本结果，但在 10.12 的证明中隐含地使用了更一般的\Term{form}{形式}；我们据此重写了本命题。}
设 $S$ 为一个\Term{scheme}{概形}，$X$ 为一个在 $S$ 上\Term{flat}{平坦}且\Term{locally-finitely-presented}{局部有限表示}的 $S$-\Term{scheme}{概形}，$H$ 为一个在 $S$ 上\Term{locally-finitely-presented}{局部有限表示}且\Term{reduced}{纤维约化}的 $S$-\Term{group}{群}，$f:X\to H$ 为一个 $S$-\Term{scheme}{概形}\Term{morphism}{态射}，$f^\infty:X^\infty\to H$ 为上面引入的 $S$-\Term{morphism}{态射}。以下条件\Term{equivalence}{等价}：

(i) $H$ 表示\Term{presheaf}{预层} $\langle\operatorname{Im}f\rangle$ 的\Term{associated-sheaf}{伴随} (fppf) $S$-\Term{sheaf}{层}。

(ii) $f^\infty$ 对 (fppf) \Term{topology}{拓扑}是\Term{covering-morphism}{覆盖态射}。

\SourcePage{391}
(iii) $f^\infty$ 为\Term{surjectivity}{满射}，即 $H=\bigcup_{n\geq1}f^n(X^n)$。

若再假设 $H$ \Term{quasi-compact}{拟紧}，则这些条件还\Term{equivalence}{等价}于下述条件：

(iv) 存在整数 $n$，使 $f^n:X^n\to H$ 对 (fppf) \Term{topology}{拓扑}是\Term{covering-morphism}{覆盖态射}（因而更是\Term{surjectivity}{满射}）。

特别地，这适用于以下情形：$X$ 为一个\Term{locally-finite-type}{局部有限型}、\Term{geometrically-reduced}{几何约化}的 $k$-\Term{scheme}{概形}，$G$ 为一个\Term{locally-finite-type}{局部有限型} $k$-\Term{group}{群}，$\phi:X\to G$ 为一个 $k$-\Term{scheme}{概形}\Term{morphism}{态射}，$H=\Gamma_G(\phi)$，且 $f:X\to H$ 为 $\phi$ 诱导的\Term{morphism}{态射}。
\end{SGAresult}
\begin{proof}
(i) 中的\Term{sheaf}{层}是 $X^\infty$ 在 $f^\infty$ 下的\Term{image-sheaf}{像层}，所以由 IV 4.4.3，$H$ 表示此\Term{sheaf}{层}\Term{equivalence}{等价}于 $f^\infty$ 为\Term{covering-morphism}{覆盖态射}；这又蕴含 $f^\infty$ 为\Term{surjectivity}{满射}。反过来，假设 $f^\infty$ \Term{surjectivity}{满射}，证明它对 (fppf) \Term{topology}{拓扑}是\Term{covering-morphism}{覆盖态射}。

设 $s\in S$，$\eta$ 为\Term{fiber}{纤维} $H_s$ 的一个\Term{maximal-point}{极大点}，取 $x\in X_s^\infty$，使 $f^\infty(x)=\eta$（因为 $f^\infty$ \Term{surjectivity}{满射}，这样的 $x$ 存在）。由于\Term{fiber}{纤维} $H_s$ \Term{reduced}{约化}，$\mathscr O_{H_s,\eta}$ 是一个域，故 $f_s^\infty$ 在 $x$ 处\Term{flat}{平坦}。由于 $X^\infty,H$ 在 $S$ 上\Term{locally-finitely-presented}{局部有限表示}且 $X^\infty$ 在 $S$ 上\Term{flat}{平坦}，由\Term{fiber}{纤维}\Term{flatness}{平坦性}判据（EGA $\mathrm{IV}_3$，11.3.10），$f^\infty$ 在 $x$ 处\Term{flat}{平坦}。因此，由\Term{lemma}{引理} 7.5，\Term{morphism}{态射}
\[
\mu\circ(f^\infty\times_S f^\infty):X^\infty\times_S X^\infty\longrightarrow H
\]
对 (fppf) \Term{topology}{拓扑}是\Term{covering-morphism}{覆盖态射}。由于 $X^n\times_S X^m$ \Term{canonical}{典范}\Term{isomorphism}{同构}于 $X^{n+m}$，且在这个\Term{isomorphism}{同构}下 $\mu\circ(f^n\times_S f^m)$ 对应于 $f^{n+m}$，显然 $\mu\circ(f^\infty\times_S f^\infty)$ 经过 $f^\infty$ \Term{resolution}{分解}，所以 $f^\infty$ 对 (fppf) \Term{topology}{拓扑}是\Term{covering-morphism}{覆盖态射}。这证明了 (iii)$\Rightarrow$(ii)，从而得到 (i)、(ii)、(iii) \Term{equivalence}{等价}。

还应注意，由于\Term{morphism}{态射} $X\to S$ 与 $H\to S$ 都\Term{locally-finitely-presented}{局部有限表示}，$f:X\to H$ 也如此（参见 EGA $\mathrm{IV}_1$，1.4.3 (v)）；又由于 $\mu:H\times_S H$ 也\Term{finitely-presented}{有限表示}（参见 VIA，0.1），故每个 $f^n:X^n\to H$ 也如此。于是，由 EGA $\mathrm{IV}_1$，1.9.5 (viii)，$f^n(X^n)$ 是 $H$ 的\Term{ind-constructible}{归纳可构造}\Term{subset}{子集}。

再假设 $H$ \Term{quasi-compact}{拟紧}（此时 $S$ 也\Term{quasi-compact}{拟紧}，因为 $H\to S$ \Term{surjectivity}{满射}）。由 EGA $\mathrm{IV}_1$，1.9.9，存在 $p$ 使 $f^p(X^p)=H$。和前面一样，利用 $H$ 的\Term{fiber}{纤维}\Term{reduced}{约化}这一事实及\Term{lemma}{引理} 7.5，可知\Term{morphism}{态射} $\mu\circ(f^p\times_S f^p):X^p\times_S X^p\to H$ 对 (fppf) \Term{topology}{拓扑}是\Term{covering-morphism}{覆盖态射}；而它等于 $f^{2p}:X^{2p}\to H$，7.6 得证。
\end{proof}

\textbf{注记 7.6.1。}显然，7.6 的各\Term{equivalence}{等价}条件蕴含所考虑的\Term{sheaf}{层} $F$ \Term{representable}{可表}。一般而言，逆命题不成立：\footnote{编者注（78）：下面详细展开了原文的论述。}例如，若 $k$ 的\Term{characteristic}{特征}为 $0$，取 $G=\mathbf G_{\mathrm a,k}$，$f:\operatorname{Spec}k\to G$ 为 $G(k)$ 的点 $1$ 给出的\Term{morphism}{态射}，则 $F$ 由\Term{constant-group}{常值} $k$-\Term{group}{群} $\mathbf Z_k$ 表示，而 $\Gamma_G(f)=G$，故\Term{monomorphism}{单态射} $\mathbf Z_k\hookrightarrow G$ 不是\Term{surjectivity}{满射}。

为简单起见，置于 7.6 的特殊情形的假设下，并假设 $F$ \Term{representable}{可表}。由 EGA $\mathrm{IV}_3$，8.14.2，$F$ 在 $k$ 上
\SourcePage{392}
\Term{locally-finitely-presented}{局部有限表示}，所以问题在于\Term{dominant}{支配}\Term{monomorphism}{单态射} $F\to\Gamma_G(f)$ 是否为\Term{isomorphism}{同构}，亦即是否为\Term{closed-immersion}{闭浸入}。由 1.4.2，当 $F$ \Term{quasi-compact}{拟紧}时答案是肯定的；由 VIA，0.5.1，当 $F$ \Term{connected}{连通}时会满足这个条件，因而特别地，当 $X$ \Term{connected}{连通}且 $f(X)$ 包含 $G$ 的\Term{identity-element}{单位元}时也是如此（7.2.1）。

\begin{SGAresult}{\Term{lemma}{引理}}{7.7}
设 $k$ 为一个\Term{algebraically-closed}{代数闭}域，$G$ 为一个\Term{locally-finite-type}{局部有限型} $k$-\Term{group}{群}，$X$ 为一个\Term{geometrically-reduced}{几何约化}、\Term{locally-finite-type}{局部有限型}的 $k$-\Term{scheme}{概形}，$f:X\to G$ 为一个 $k$-\Term{morphism}{态射}，$H$ 为 $G$ 的一个\Term{group-subscheme}{群子概形}，使 $H\subset f(X)$。令
\[
\Gamma'=\bigcup_{n\geq1}f^n(X^n),\qquad
\Gamma'_0=\Gamma'\cap G(k),\qquad H_0=H(k),
\]
并假设 $H_0$ 在 $\Gamma'_0$ 中的\Term{index-subgroup}{指数}\Term{finite}{有限}。

则存在整数 $m$，使 $f^m(X^m)=\Gamma_G(f)$（参见 7.6），且 $\Gamma_G(f)$ 为 $H$ 的\Term{finite}{有限}个\Term{translation}{平移}的并。
\end{SGAresult}
\begin{proof}
对每个 $n\geq1$，$f^n(X^n)$ 为 $G$ 的\Term{ind-constructible}{归纳可构造}\Term{subset}{子集}（EGA $\mathrm{IV}_1$，1.9.5 (viii)），$\Gamma'$ 也如此。由于 $G$ 是一个 Jacobson \Term{scheme}{概形}，$\Gamma'_0$ 在 $\Gamma'$ 中\Term{dense}{稠密}。由假设，存在 $\Gamma'_0$ 中的\Term{finite}{有限}列点 $a_1,\ldots,a_r$，使 $\Gamma'_0=a_1H_0\cup\cdots\cup a_rH_0$，故
\[
\begin{aligned}
\Gamma_G(f)&=\overline{\Gamma'}=\overline{\Gamma'_0}
=\overline{a_1H_0\cup\cdots\cup a_rH_0}\\
&=\overline{a_1H_0}\cup\cdots\cup\overline{a_rH_0}
=a_1\overline{H_0}\cup\cdots\cup a_r\overline{H_0},
\end{aligned}
\]
最后一个等式是因为由 $a_i$ 给出的\Term{translation}{平移}为 $G$ 到自身的一个\Term{homeomorphism}{同胚}。所以 $\Gamma_G(f)=a_1H\cup\cdots\cup a_rH$。另一方面，显然存在整数 $p$，使每个 $a_i$（$1\leq i\leq r$）都属于 $f^p(X^p)$。最后，由于 $H\subset f(X)$，对每个 $i$ 有 $a_iH\subset f^{p+1}(X^{p+1})$，故 $f^{p+1}(X^{p+1})=\Gamma_G(f)$。
\end{proof}

\begin{SGAresult}{命题}{7.8}
设 $G$ 为一个\Term{locally-finite-type}{局部有限型} $k$-\Term{group}{群}，$A,B$ 为 $G$ 的两个\Term{geometrically-reduced}{几何约化}的 $k$-\Term{group-subscheme}{群子概形}（所以在其\Term{irreducible-component}{不可约分支}的\Term{generic-point}{泛点}处\Term{smooth}{光滑}，进而由 1.3 知处处\Term{smooth}{光滑}）。假设满足下述条件 a)、b) 之一：

a) $A,B$ 均\Term{normal-subgroup}{正规}且在 $k$ 上\Term{finite-type}{有限型}。

b) $A$ \Term{connected}{连通}，且 $B$ 在 $k$ 上\Term{finite-type}{有限型}。

则 $(A,B)$ 在 $k$ 上\Term{finite-type}{有限型}，并表示 $A,B$ 在 $G$ 中的\Term{commutator-subgroup}{交换子群}\Term{presheaf}{预层}对于 (fppf)（或 (fpqc)）\Term{topology}{拓扑}的\Term{associated-sheaf}{伴随层}。此外，\footnote{编者注（79）：对下文作了更精确的表述，并在证明中详细展开了归结到 $k$ 为\Term{algebraically-closed}{代数闭}的情形的步骤。}$k$-\Term{group}{群} $(A,B^0)$ 与 $(A^0,B)$ \Term{connected}{连通}，且
\[
(A,B)^0=(A,B^0)\cdot(A^0,B).
\]
\end{SGAresult}
\begin{proof}
由 7.6，为证明 $(A,B)$ 是所需的\Term{associated-sheaf}{伴随层}，只须证明存在整数 $n$，使 $\nu^n((A\times_k B)^n)=(A,B)$（采用 7.2.2 的记号）。为证明此式及其余两条断言，可以假设
\SourcePage{393}
$k$ \Term{algebraically-closed}{代数闭}。事实上，设 $\overline k$ 为 $k$ 的一个\Term{algebraic-closure}{代数闭包}。由 7.1 (iii) 与 VIA 2.4，采用显然的记号，有
\[
\begin{gathered}
(A,B)_{\overline k}=(A_{\overline k},B_{\overline k}),\qquad
((A,B)^0)_{\overline k}=(A_{\overline k},B_{\overline k})^0,\\
(A,B^0)_{\overline k}=(A_{\overline k},B^0_{\overline k}),\qquad\text{等等。}
\end{gathered}
\]
所以，若证明 $(A_{\overline k},B_{\overline k})$ 在 $\overline k$ 上\Term{finite-type}{有限型}（相应地，$(A_{\overline k},B^0_{\overline k})$ 与 $(A^0_{\overline k},B_{\overline k})$ \Term{connected}{连通}，且\Term{morphism}{态射} $(A_{\overline k},B^0_{\overline k})\times_{\overline k}(A_{\overline k},B^0_{\overline k})\to(A_{\overline k},B_{\overline k})^0$ \Term{surjectivity}{满射}），则由 EGA $\mathrm{IV}_2$，2.7.1 与 2.6.1，类似的断言在 $k$ 上也成立。

设 $B^1,\ldots,B^p$ 为 $B$ 除\Term{neutral-component}{单位连通分支} $B^0$ 以外的\Term{connected-component}{连通分支}（由于 $B$ 假设在 $k$ 上\Term{finite-type}{有限型}，因而为\Term{noetherian}{Noether}，它们只有\Term{finite}{有限}个）；在 (a) 中，同样记 $A^1,\ldots,A^q$ 为 $A$ 的这些\Term{connected-component}{连通分支}。（在 (b) 中，只考虑 $A^0=A$。）记 $\nu_{ij}$ 为 $\nu$ 在 $A^i\times_k B^j$ 上的\Term{restriction}{限制}。每个 $A^i,B^j$ 都\Term{irreducible}{不可约}（VIA 2.4.1），故 $A^0\times_k B^j$ 与 $A^i\times_k B^0$ 也如此。由于 $G$ 的\Term{identity-element}{单位元}属于 $A^0,B^0$，它属于 $\nu_{0j}(A^0\times_k B^j)$ 与 $\nu_{i0}(A^i\times_k B^0)$。于是每个 $\Gamma_G(\nu_{0j})$ 与 $\Gamma_G(\nu_{i0})$ 都\Term{connected}{连通}（7.2.1）。同样，若 $u_{0j}$（相应地，$u_{i0}$）表示 $\Gamma_G(\nu_{0j})$（相应地，$\Gamma_G(\nu_{i0})$）到 $G$ 的\Term{injectivity}{单射}，则
\[
(A^0,B)=\Gamma_G((u_{0j})_{j=0}^p),\qquad
(A,B^0)=\Gamma_G((u_{i0})_{i=0}^q)
\]
都是\Term{connected}{连通}的。此外，由 7.4 及上述构造容易\Term{pushout}{推出}，存在指标 $r$，使 $(A^0,B)$ 与 $(A,B^0)$ 都包含在 $\nu^r((A\times_k B)^r)$ 中。在 b) 的情形，$(A,B)=(A^0,B)$，证明完毕。

现在考虑 (a)。\footnote{编者注（80）：下面利用增补 7.1.1，详细展开了原文的论证。}已知 $(A^0,B)$ 与 $(A,B^0)$ 都是 $G$ 的\Term{smooth}{光滑}\Term{connected}{连通} $k$-\Term{subgroup}{子群}，因而\Term{finite-type}{有限型}（参见 VIA，2.4）。另一方面，由于 $A^0$ 是 $A$ 的\Term{characteristic-subgroup}{特征}\Term{subgroup}{子群}（参见 VIA，2.6.5），它是 $G$ 的\Term{normal-subgroup}{正规}\Term{subgroup}{子群}；由 7.3 (v)，$(A^0,B)$ 是 $G$ 的\Term{normal-subgroup}{正规}\Term{subgroup}{子群}，$(A,B^0)$ 亦然。因此，由 7.1.1，$G$ 中由 $(A,B^0)$ 与 $(A^0,B)$ \Term{generated-subgroup}{生成的子群} $H$ 正是 $(A,B^0)\cdot(A^0,B)$。特别地，有 $H\subset\nu^{2r}((A\times_k B)^{2r})$。

对 $G$ 中任意在\Term{group-law}{群运算}下\Term{stable}{稳定}的\Term{subset}{子集} $X$（参见 3.0），用 $X_0$ 表示属于 $X$ 的 $G$ 的 $k$-点群。令 $\Gamma'=\bigcup_{q\geq1}\nu^q((A\times_k B)^q)$。由下面的命题 7.9，有
\[
(A^0,B)_0=(A^0_0,B_0),\qquad
(A,B^0)_0=(A_0,B^0_0),\qquad
\Gamma'_0=(A_0,B_0).
\]
于是 $H_0=(A^0_0,B_0)\cdot(A_0,B^0_0)$ 在 $\Gamma'_0$ 中的\Term{index-subgroup}{指数}\Term{finite}{有限}（Bible，第 3 讲，附录），因为 $A_0,B_0$ \Term{normal-subgroup}{正规}，且 $A^0_0$（相应地，$B^0_0$）在 $A_0$（相应地，$B_0$）中的\Term{index-subgroup}{指数}\Term{finite}{有限}。因此满足\Term{lemma}{引理} 7.7 的条件：由于 $H\subset\nu^{2r}((A\times_k B)^{2r})$，存在整数 $m$，使 $\nu^{2rm}((A\times_k B)^{2rm})=(A,B)$，并存在 $G$ 的\Term{finite}{有限}列 $k$-点 $a_1,\ldots,a_n$，使 $(A,B)=a_1H\cup\cdots\cup a_nH$。由于 $H$ 在 $k$ 上\Term{finite-type}{有限型}，每个 $a_iH$ 都\Term{quasi-compact}{拟紧}，故它们的并 $(A,B)$ \Term{quasi-compact}{拟紧}，从而在 $k$ 上\Term{finite-type}{有限型}。由于 $H$ \Term{irreducible}{不可约}，每个 $a_iH$ 也\Term{irreducible}{不可约}；又因 $e\in H$，显然 $H=(A,B)^0$。
\end{proof}

\SourcePage{394}
\begin{SGAresult}{命题}{7.9}
设 $k$ 为一个\Term{algebraically-closed}{代数闭}域，$G$ 为一个\Term{locally-finite-type}{局部有限型} $k$-\Term{group}{群}。

(i) 设 $A,B$ 为 $G$ 的两个\Term{ind-constructible}{归纳可构造}\Term{subset}{子集}。用 $A_0$ 表示属于 $A$ 的 $G$ 的\Term{rational-point}{有理点}的\Term{set}{集合}。则 $(A\cdot B)_0=A_0\cdot B_0$，右边的乘积在群 $G(k)$ 中进行。

(ii) 设 $X$ 为一个\Term{geometrically-reduced}{几何约化}且\Term{locally-finite-type}{局部有限型}的 $k$-\Term{scheme}{概形}，$f:X\to G$ 为一个 $k$-\Term{morphism}{态射}。令 $\Gamma'_G(f)=\bigcup_{n\geq1}f^n(X^n)$。则 $\Gamma'_G(f)_0$ 是 $G(k)$ 中由 $f(X)_0$ \Term{generated-subgroup}{生成的子群}。

(iii) 特别地，设 $A,B$ 为 $G$ 的两个\Term{smooth}{光滑}\Term{group-subscheme}{群子概形}，令 $\Gamma'=\bigcup_{n\geq1}\nu^n(A\times_k B)^n$（采用 7.2.2 的记号）。则 $\Gamma'_0$ 是 $A(k),B(k)$ 在 $G(k)$ 中的\Term{commutator-subgroup}{交换子群}。
\end{SGAresult}
\begin{proof}
先证 (i)。显然 $A_0\cdot B_0\subset(A\cdot B)_0$。反过来，设 $z\in(A\cdot B)_0$。则 $\mu^{-1}(z)$ 是 $G\times_k G$ 的一个闭集，而 $A\times_k B$（参见 3.0）是 $G\times_k G$ 的\Term{ind-constructible}{归纳可构造}\Term{subset}{子集}，所以 $\mu^{-1}(z)\cap(A\times_k B)$ 是 $G\times_k G$ 的一个\Term{nonempty}{非空}\Term{ind-constructible}{归纳可构造}\Term{subset}{子集}。由 EGA $\mathrm{IV}_3$，10.4.8，它包含 $G\times_k G$ 的一个\Term{rational-point}{有理点}，其\Term{projection}{投影} $x,y$ 为 $G$ 的\Term{rational-point}{有理点}，满足 $x\in A$、$y\in B$ 及 $x\cdot y=z$，所以 $(A,B)_0=A_0\cdot B_0$。

为证明 (ii)，注意，由于 $f^n$ \Term{locally-finite-type}{局部有限型}，$f^n(X^n)$ 是 $G$ 的\Term{ind-constructible}{归纳可构造}\Term{subset}{子集}（EGA $\mathrm{IV}_4$，1.9.5 (viii)）。令 $A=f(X)_0$，由 (i) 可归纳证明 $f^n(X^n)_0=(A\cup A^{-1})^n$，因此
\[
\Gamma'_G(f)_0=\bigcup_{n\geq1}f^n(X^n)_0
=\bigcup_{n\geq1}(A\cup A^{-1})^n,
\]
它正是 $G(k)$ 中由 $A=f(X)_0$ \Term{generated-subgroup}{生成的子群}。

最后，(iii) 由 (ii) 及定义得出。
\end{proof}

\begin{SGAresult}{推论}{7.10}
在 7.8 的条件下，若 $k$ \Term{algebraically-closed}{代数闭}，则 $(A,B)(k)$ 是 $A(k),B(k)$ 在 $G(k)$ 中的\Term{commutator-subgroup}{交换子群}。
\end{SGAresult}
事实上，只须应用 7.9 (iii)、7.8 和 7.6。

\SGAsection{8}{可解或幂零的群概形}
\textbf{8.1。}设 $\mathscr C$ 为一个带有\Term{topology}{拓扑} $\mathscr T$ 的\Term{category}{范畴}（参见 IV 第 4 节）。对 $\mathscr C$ 上的每个\Term{presheaf}{预层} $P$，记 $P^\flat$ 为它的\Term{associated-sheaf}{伴随层}。

设 $G$ 为 $\mathscr C$ 上的一个\Term{group-presheaf}{群预层}，$A,B$ 为 $G$ 的两个\Term{group-subpresheaf}{群子预层}，并令 $\operatorname{Comm}(A,B)$ 为 $A,B$ 在 $G$ 中的\Term{commutator-subgroup}{交换子群}\Term{presheaf}{预层}；即对每个 $S\in\operatorname{Ob}\mathscr C$，$\operatorname{Comm}(A,B)(S)$ 是 $G(S)$ 中由\Term{commutator}{交换子} $aba^{-1}b^{-1}$（$a\in A(S),b\in B(S)$）\Term{generated-subgroup}{生成的子群}。记
\[
\operatorname{Comm}_{\mathscr T}(A,B)=\operatorname{Comm}(A,B)^\flat.
\]

\SourcePage{395}
在 8.2 的证明中，需要以下结果。\footnote{编者注（81）：原文仅提及这些结果，未给证明；我们将它们单列为\Term{lemma}{引理} 8.1.1 与 8.1.2，并详细展开了证明。}
\begin{SGAresult}{\Term{lemma}{引理}}{8.1.1}
设 $A\subset G$ 为\Term{group-sheaf}{群层}，且 $A$ 在 $G$ 中\Term{normal-subgroup}{正规}。

(i) $\operatorname{Comm}_{\mathscr T}(G,A)$ 是 $G$ 的满足如下条件的最小\Term{normal-subgroup}{正规}\Term{group-subsheaf}{群子层} $C$：商\Term{presheaf}{预层} $A/C$ 的\Term{associated-sheaf}{伴随层} $(A/C)^\flat$ 在 $(G/C)^\flat$ 中\Term{central-subgroup}{中心}。

(ii) 特别地，$\operatorname{Comm}_{\mathscr T}(G,G)$ 是 $G$ 的使商\Term{sheaf}{层} $(G/C)^\flat$ \Term{commutative-group}{交换}的最小\Term{normal-subgroup}{正规}\Term{group-subsheaf}{群子层} $C$。
\end{SGAresult}
\begin{proof}
显然 (ii) 是 (i) 中 $A=G$ 的特例，所以只须证明 (i)。设 $C$ 为 $A$ 的一个在 $G$ 中\Term{normal-subgroup}{正规}的\Term{group-subsheaf}{群子层}，且商\Term{sheaf}{层} $(A/C)^\flat$ 在 $(G/C)^\flat$ 中\Term{central-subgroup}{中心}。由\Term{lemma}{引理} IV 4.4.8.1，\Term{presheaf}{预层} $A/C,G/C$ \Term{separated-presheaf}{分离}，因而由 IV 4.3.11，下图的所有\Term{morphism}{态射}都是\Term{monomorphism}{单态射}：
\[
\begin{tikzcd}
A/C\arrow[r,hook]\arrow[d,hook] & G/C\arrow[d,hook]\\
(A/C)^\flat\arrow[r,hook] &(G/C)^\flat.
\end{tikzcd}
\]
由于 $(A/C)^\flat$ 在 $(G/C)^\flat$ 中\Term{central-subgroup}{中心}，$A/C$ 在 $G/C$ 中\Term{central-subgroup}{中心}，故 $\operatorname{Comm}(G,A)\subset C$，所以由 IV 4.3.12，$C$ 包含 $\operatorname{Comm}_{\mathscr T}(G,A)$。

反过来，$\operatorname{Comm}(G,A)$ 是 $A$ 的一个在 $G$ 中\Term{normal-subgroup}{正规}的\Term{group-subpresheaf}{群子预层}，并且\Term{separated-presheaf}{分离}（参见 IV 4.3.1，编者注（24））。所以由 IV 4.4.8.2 (i) 及 IV 4.3.11，$C=\operatorname{Comm}_{\mathscr T}(G,A)$ 是 $A$ 的一个在 $G$ 中\Term{normal-subgroup}{正规}且包含 $\operatorname{Comm}(G,A)$ 的\Term{group-subsheaf}{群子层}。因此\Term{presheaf}{预层} $A/C$ 在 $G/C$ 中\Term{central-subgroup}{中心}；由 IV 4.4.8.2 (ii)，$(A/C)^\flat$ 在 $(G/C)^\flat$ 中\Term{central-subgroup}{中心}。\Term{lemma}{引理} 8.1.1 得证。
\end{proof}

\begin{SGAresult}{\Term{lemma}{引理}}{8.1.2}
设 $G$ 为一个\Term{group-sheaf}{群层}，$A,B$ 为 $G$ 的两个\Term{group-subpresheaf}{群子预层}。

(i) \Term{morphism}{态射} $\tau:\operatorname{Comm}(A,B)\to\operatorname{Comm}(A^\flat,B^\flat)$ 是一个\Term{covering-morphism}{覆盖}\Term{monomorphism}{单态射}。

(ii) 因此有\Term{isomorphism}{同构}
\[
\operatorname{Comm}_{\mathscr T}(A,B)\xrightarrow{\sim}
\operatorname{Comm}_{\mathscr T}(A^\flat,B^\flat).
\]
\end{SGAresult}
\begin{proof}
(i) 由于 $A$（相应地，$B$）是 $G$ 的一个\Term{subpresheaf}{子预层}，$A^\flat$（相应地，$B^\flat$）是 $G$ 中包含 $A$（相应地，$B$）的一个\Term{subpresheaf}{子预层}，因此 $\tau$ 是\Term{monomorphism}{单态射}。

证明 $\tau$ 为\Term{covering-morphism}{覆盖态射}。设 $S\in\operatorname{Ob}\mathscr C$，$g\in\operatorname{Comm}(A^\flat,B^\flat)(S)$。则存在整数 $n\geq1$，并对 $i=1,\ldots,n$ 存在元素 $a'_i\in A^\flat(S)$、$b'_i\in B^\flat(S)$ 及 $\varepsilon_i\in\{\pm1\}$，使
\[
g=(a'_1,b'_1)^{\varepsilon_1}\cdots(a'_n,b'_n)^{\varepsilon_n},
\]
\SourcePage{396}
其中 $(a,b)$ 表示\Term{commutator}{交换子} $aba^{-1}b^{-1}$，并且存在 $S$ 的一个\Term{refinement}{加细} $R$，使对所有 $i$ 有 $a'_i\in A(R)$、$b'_i\in B(R)$。于是 $g$ 是如下复合\Term{morphism}{态射}：
\[
R\xrightarrow{\,(a'_1,\ldots,b'_n)\,}(A\times B)^n
\xrightarrow{\,\Phi^{\varepsilon_1,\ldots,\varepsilon_n}\,}\operatorname{Comm}(A,B),
\]
其中 $\Phi=\Phi^{\varepsilon_1,\ldots,\varepsilon_n}$ 为对任意 $T\in\operatorname{Ob}\mathscr C$ 和 $a_i\in A(T),b_i\in B(T)$ 在\Term{set}{集合}意义下由
\[
\Phi(a_1,b_1,\ldots,a_n,b_n)=(a_1,b_1)^{\varepsilon_1}\cdots(a_n,b_n)^{\varepsilon_n}
\]
定义的\Term{morphism}{态射}。这表明 $g\in\operatorname{Comm}(A,B)(R)$，如 IV 4.3.11 (i) 的证明一样，可知 $\operatorname{Comm}(A,B)\to\operatorname{Comm}(A^\flat,B^\flat)$ 为\Term{covering-morphism}{覆盖态射}。

由于 $\operatorname{Comm}(A^\flat,B^\flat)\to\operatorname{Comm}(A^\flat,B^\flat)^\flat$ 也是\Term{covering-morphism}{覆盖}\Term{monomorphism}{单态射}（IV 4.3.11 (iv)），$\operatorname{Comm}(A,B)\to\operatorname{Comm}(A^\flat,B^\flat)^\flat$ 也如此；因此由 IV 4.3.12，得到\Term{isomorphism}{同构}
\[
\operatorname{Comm}(A,B)^\flat\xrightarrow{\sim}
\operatorname{Comm}(A^\flat,B^\flat)^\flat.
\]
\Term{lemma}{引理} 8.1.2 得证。
\end{proof}

\begin{SGAresult}{命题}{8.2}
设 $\mathscr C$ 为一个\Term{category}{范畴}，$\mathscr T$ 为 $\mathscr C$ 上的一个\Term{topology}{拓扑}，$G$ 为 $\mathscr C$ 上的一个\Term{group-sheaf}{群层}，$n$ 为一个非负整数。下列条件\Term{equivalence}{等价}：

(i) 令 $K_0=G$，并对 $p\geq1$ 令 $K_p=\operatorname{Comm}(K_{p-1},K_{p-1})$（相应地，$K_p=\operatorname{Comm}(G,K_{p-1})$），则 $K_n$ 为单位\Term{group-presheaf}{群预层}。

(ii) 令 $K'_0=G$，并对 $p\geq1$ 令 $K'_p=\operatorname{Comm}_{\mathscr T}(K'_{p-1},K'_{p-1})$（相应地，$K'_p=\operatorname{Comm}_{\mathscr T}(G,K'_{p-1})$），则 $K'_n$ 为单位\Term{group-sheaf}{群层}。

(iii) 存在 $G$ 的\Term{normal-subgroup}{正规}\Term{subsheaf}{子层}序列 $G=H_0\supset H_1\supset\cdots\supset H_n$，使对每个 $i$，商\Term{sheaf}{层} $H_i/H_{i+1}$ \Term{commutative-group}{交换}（相应地，在 $G/H_{i+1}$ 中\Term{central-subgroup}{中心}），且 $H_n$ 为单位\Term{sheaf}{层}。
\end{SGAresult}
\begin{proof}
显然 $K_n\subset K'_n$，所以 (ii) 蕴含 (i)。证明 (i) 蕴含 (ii)。有 $K'_1=\operatorname{Comm}_{\mathscr T}(G,G)=K_1^\flat$，再由\Term{lemma}{引理} 8.1.2 归纳可得，对任意 $p$，$K'_p=K_p^\flat$。因此，若 $K_n$ 为单位\Term{presheaf}{预层}，则 $K'_n=K_n^\flat$ 为单位\Term{sheaf}{层}。

最后，由\Term{lemma}{引理} 8.1.1，条件 (ii) 与 (iii) \Term{equivalence}{等价}。
\end{proof}

\begin{SGAresult}{定义}{8.2.1}
当这些条件满足时，称\Term{sheaf}{层} $G$ 为 $n$ 类\Term{solvable}{可解}（相应地，$n$ 类\Term{nilpotent}{幂零}）的。若存在整数 $n$ 使这些条件满足，则称 $G$ \Term{solvable}{可解}（相应地，\Term{nilpotent}{幂零}）。
\end{SGAresult}
注意，由条件 (i)，这些概念不依赖于\Term{topology}{拓扑} $\mathscr T$。

\begin{SGAresult}{命题}{8.3}
设 $k$ 为一个域，$S$ 为一个\Term{nonempty}{非空} $k$-\Term{scheme}{概形}，$\Omega$ 为 $k$ 的一个\Term{algebraically-closed}{代数闭}\Term{field-extension}{扩域}，$G$ 为一个\Term{smooth}{光滑}、\Term{finite-type}{有限型} $k$-\Term{group}{群}。下列条件\Term{equivalence}{等价}：

(i) $G$ 为 $n$ 类\Term{solvable}{可解}（相应地，$n$ 类\Term{nilpotent}{幂零}）。

(ii) $G\times_k S$ 为 $n$ 类\Term{solvable}{可解}（相应地，$n$ 类\Term{nilpotent}{幂零}）。

(iii) 群 $G(\Omega)$ 为 $n$ 类\Term{solvable}{可解}（相应地，$n$ 类\Term{nilpotent}{幂零}）。

\SourcePage{397}
(iv) 令 $K_0=G$，并对 $p\geq1$ 考虑 $k$-\Term{group}{群} $K_p=(K_{p-1},K_{p-1})$（相应地，$K_p=(G,K_{p-1})$）（参见 7.2.2），则 $K_n$ 为单位 $k$-\Term{group}{群}。
\end{SGAresult}
\begin{proof}
由于\Term{commutator-subgroup}{交换子群}\Term{presheaf}{预层}的形成与\Term{base-change}{基变换}交换（IV 4.1.3），由命题 8.2，条件 (i) 与 (ii) \Term{equivalence}{等价}。

由 7.8，$k$-\Term{group}{群} $K_p=(K_{p-1},K_{p-1})$（相应地，$K_p=(G,K_{p-1})$）表示\Term{sheaf}{层} $\operatorname{Comm}_{\mathscr T}(K_{p-1},K_{p-1})$（相应地，$\operatorname{Comm}_{\mathscr T}(G,K_{p-1})$），其中 $\mathscr T$ 为 (fppf)（或 (fpqc)）\Term{topology}{拓扑}；因此 (i) 与 (iv) \Term{equivalence}{等价}。

为证明 (iii) 与 (iv) \Term{equivalence}{等价}，可假设 $k=\Omega$，此时它们的\Term{equivalence}{等价}性由 7.10 得出。
\end{proof}

\begin{SGAresult}{命题}{8.4}
设 $G$ 为一个\Term{finitely-presented}{有限表示} $S$-\Term{group}{群}，且对每个 $s\in S$，$G_s$ 在 $\kappa(s)$ 上\Term{smooth}{光滑}。设 $T$ 为使 $G_s$ \Term{solvable}{可解}（相应地，\Term{nilpotent}{幂零}）的 $s\in S$ 组成的\Term{set}{集合}。

(i) 则 $T$ 在 $S$ 中\Term{locally-constructible}{局部可构造}。

(ii) 若再假设 $G$ 在 $S$ 上\Term{flat}{平坦}且\Term{separated}{分离}（即 $G$ 在 $S$ 上\Term{smooth}{光滑}、\Term{quasi-compact}{拟紧}且\Term{separated}{分离}），则 $T$ 在 $S$ 中闭。
\end{SGAresult}
\begin{proof}
显然可假设 $S$ 为环 $A$ 的\Term{affine}{仿射}\Term{scheme}{概形}。由 10.1 和 10.10 b)，\footnote{在本证明中将使用第 10 节建立的结果；第 10 节与第 9 节一样，不依赖于本节（第 8 节）。}存在 $A$ 的一个\Term{noetherian}{Noether}\Term{subring}{子环} $A'$ 及一个\Term{finite-type}{有限型} $A'$-\Term{group}{群} $G'$，使 $G'\otimes_{A'}A$ \Term{isomorphism}{同构}于 $G$。由 EGA $\mathrm{IV}_3$，11.2.6 与 8.10.5，\footnote{编者注（82）：将 10.8.5 更正为 8.10.5。}若 $G$ 在 $S$ 上\Term{flat}{平坦}且\Term{separated}{分离}，可假设 $G'$ 在 $S'=\operatorname{Spec}A'$ 上\Term{flat}{平坦}且\Term{separated}{分离}。\footnote{编者注（83）：下文作了详细展开。}由于 $G$ 在 $S$ 上\Term{finitely-presented}{有限表示}，使 $G_s$ \Term{geometrically-reduced}{几何约化}（\Term{equivalence}{等价}地，在 $\kappa(s)$ 上\Term{smooth}{光滑}）的 $s\in S$ 组成的\Term{set}{集合}\Term{locally-constructible}{局部可构造}（EGA $\mathrm{IV}_3$，9.7.7）。故由 EGA $\mathrm{IV}_3$，9.3.3，可假设对每个 $s'\in S'$，$G'_{s'}$ 都在 $\kappa(s')$ 上\Term{smooth}{光滑}。另一方面，若 $s'$ 表示 $s$ 在 $S'$ 中的像，则有 $G'_{s'}\otimes_{\kappa(s')}\kappa(s)\simeq G_s$。所以由 8.3，$G_s$ \Term{solvable}{可解}（相应地，\Term{nilpotent}{幂零}），当且仅当 $G'_{s'}$ 具有相同性质。因此归结到了 $S$ 为\Term{affine}{仿射}\Term{noetherian}{Noether}\Term{scheme}{概形}的情形。

现在证明 $T$ \Term{constructible}{可构造}。应用判据（EGA $0_{\mathrm{III}}$，9.2.3），并如上论证，可知只须证明：当 $S$ 为\Term{noetherian}{Noether}且\Term{integral}{整}时，$T$ 或 $S-T$ 包含 $S$ 的一个\Term{nonempty}{非空}\Term{open-subset}{开集}。

因此假设 $S$ \Term{integral}{整}且\Term{noetherian}{Noether}，\Term{generic-point}{泛点}为 $\eta$。对任意 $s\in S$，令 $K_s^0=G_s$，并令 $K_s^p=(K_s^{p-1},K_s^{p-1})$（相应地，$K_s^p=(G,K_s^{p-1})$）。先证明\Term{closed-subscheme}{闭子概形}序列 $K_\eta^p$ 最终\Term{stable}{稳定}。由 7.3 (v)，每个 $K_\eta^p$ 都\Term{normal-subgroup}{正规}，故 $K_\eta^p$ 构成递减序列；由于 $G_\eta$ 为\Term{noetherian}{Noether}，该序列最终\Term{stable}{稳定}。因此存在整数 $n$，使对每个 $p\geq n$，$K_\eta^p=K_\eta^n$。

另一方面，由 10.12.1 和 10.13，存在 $S$ 的一个\Term{nonempty}{非空}\Term{open-subset}{开集} $S'$ 及
\SourcePage{398}
一个\Term{finitely-presented}{有限表示} $S'$-\Term{group}{群} $D$，使对每个 $s\in S'$，$D_s=K_s^n$ 且 $(D_s,D_s)=D_s$（相应地，$(G_s,D_s)=D_s$）。可假设 $S'=S$。于是对任意 $s\in S$ 及任意 $p\geq n$，$D_s=K_s^p$，所以 $G_s$ \Term{solvable}{可解}（相应地，\Term{nilpotent}{幂零}），当且仅当 $D_s$ \Term{isomorphism}{同构}于单位 $\kappa(s)$-\Term{group}{群}。

由 EGA $\mathrm{IV}_3$，9.6.1 (xi)，使\Term{structure-morphism}{结构态射} $D_s\to\operatorname{Spec}\kappa(s)$ 为\Term{isomorphism}{同构}的 $s\in S$ 组成的\Term{set}{集合}\Term{constructible}{可构造}，\footnote{编者注（84）：这里利用了 $S$ 假设为\Term{affine}{仿射}，因而\Term{quasi-compact}{拟紧}且\Term{quasi-separated}{拟分离}（参见 EGA $0_{\mathrm{III}}$，9.1.12）。}故或者是\Term{nowhere-dense}{无处稠密}的，或者包含 $S$ 的一个\Term{nonempty}{非空}\Term{open-subset}{开集}。这样就证明了 $T$ \Term{locally-constructible}{局部可构造}。

再证明，若 $G$ 在 $S$ 上\Term{flat}{平坦}且\Term{separated}{分离}，则 $T$ 闭。由于 $T$ \Term{locally-constructible}{局部可构造}，$T$ 闭的充要条件是它在\Term{specialization}{特殊化}下\Term{stable}{稳定}（参见 EGA $\mathrm{IV}_1$，1.10.1）。

设 $s\in S$，$s'$ 为 $s$ 在 $S$ 中的一个\Term{specialization}{特殊化}。由于已归结到 $S$ 为\Term{noetherian}{Noether}的情形，由 EGA II，7.1.9，存在一个\Term{discrete-valuation-ring}{离散赋值环} $A$ 及一个\Term{morphism}{态射} $\operatorname{Spec}(A)\to S$，使 $s$（相应地，$s'$）为 $\operatorname{Spec}(A)$ 的\Term{generic-point}{泛点} $\alpha$（相应地，\Term{closed-point}{闭点} $a$）的像。于是只须证明，令 $G'=G\otimes_S A$，若 $G'_\alpha$ \Term{solvable}{可解}（相应地，\Term{nilpotent}{幂零}），则 $G'_a$ 也如此。注意，由于 $G$ 在 $S$ 上\Term{flat}{平坦}且\Term{separated}{分离}，$G'_\alpha$ 在 $A$ 上\Term{flat}{平坦}且\Term{separated}{分离}，所以归结到 $S$ 是一个\Term{discrete-valuation-ring}{离散赋值环} $A$ 的谱的情形。

由于假设 $G_\alpha$ \Term{solvable}{可解}（相应地，\Term{nilpotent}{幂零}），存在整数 $q$，使 $K_\alpha^q$（采用前面引入的记号）\Term{isomorphism}{同构}于单位 $\kappa(\alpha)$-\Term{group}{群}。对每个 $n$，以 $\overline{K_\alpha^n}$ 表示 $K_\alpha^n$ 在 $G$ 中的\Term{schematic-closure}{概形闭包}（EGA $\mathrm{IV}_2$，2.8.5 的意义下）。对 $p$ 作归纳，证明 $(\overline{K_\alpha^p})_a\supset K_a^p$。先注意，由于 $G$ 在 $A$ 上\Term{flat}{平坦}，$\overline{G_\alpha}=G$（EGA $\mathrm{IV}_2$，2.8.5），故 $(\overline{K_\alpha^0})_a=K_a^0$。

设 $p\geq0$。假设已证明 $K_a^p\subset(\overline{K_\alpha^p})_a$，并以 $\nu_a,\nu_\alpha,\overline\nu,\overline\nu_a$ 表示以下按 7.2.2 定义的\Term{morphism}{态射}：
\[
\begin{array}{c|c|c}
&\text{可解情形}&\text{幂零情形}\\ \hline
\nu_a&K_a^p\times_{\kappa(a)}K_a^p\to G_a
&G_a\times_{\kappa(a)}K_a^p\to G_a\\
\nu_\alpha&K_\alpha^p\times_{\kappa(\alpha)}K_\alpha^p\to G_\alpha
&G_\alpha\times_{\kappa(\alpha)}K_\alpha^p\to G_\alpha\\
\overline\nu&\overline{K_\alpha^p}\times_A\overline{K_\alpha^p}\to G
&G\times_A\overline{K_\alpha^p}\to G\\
\overline\nu_a&(\overline{K_\alpha^p})_a\times_{\kappa(a)}(\overline{K_\alpha^p})_a\to G_a
&G_a\times_{\kappa(a)}(\overline{K_\alpha^p})_a\to G_a.
\end{array}
\]
由于 $\nu_\alpha$ 经过 $K_\alpha^{p+1}$ \Term{resolution}{分解}，$\overline\nu$ 便经过 $\overline{K_\alpha^{p+1}}$ \Term{resolution}{分解}；后者显然是 $G$ 的一个\Term{group-subscheme}{群子概形}，故 $\overline{K_\alpha^{p+1}}$ 包含 $\Gamma_G(\overline\nu)$。由 7.1 (iii)，$\Gamma_G(\overline\nu)_a=\Gamma_{G_a}(\overline\nu_a)$；由归纳假设，有
\[
\begin{gathered}
K_a^p\times_{\kappa(a)}K_a^p
\subset(\overline{K_\alpha^p})_a\times_{\kappa(a)}(\overline{K_\alpha^p})_a,\\
\text{相应地，}\quad
G_a\times_{\kappa(a)}K_a^p\subset G_a\times_{\kappa(a)}(\overline{K_\alpha^p})_a,
\end{gathered}
\]
故
\[
K_a^{p+1}=\Gamma_{G_a}(\nu_a)\subset\Gamma_{G_a}(\overline\nu_a)
=\Gamma_G(\overline\nu)_a\subset(\overline{K_\alpha^{p+1}})_a.
\]
\SourcePage{399}
又因为 $K_\alpha^q$ \Term{isomorphism}{同构}于单位 $\kappa(\alpha)$-\Term{group}{群}，而单位 $A$-\Term{group}{群}在 $A$ 上\Term{flat}{平坦}且\Term{isomorphism}{同构}于 $G$ 的一个\Term{closed-subscheme}{闭子概形}（因为 $G$ 在 $A$ 上\Term{separated}{分离}，参见 5.1），所以由 EGA $\mathrm{IV}_2$，2.8.5，\Term{schematic-closure}{概形闭包} $\overline{K_\alpha^q}$ \Term{isomorphism}{同构}于单位 $A$-\Term{group}{群}。刚才已证明 $K_a^q\subset(\overline{K_\alpha^q})_a$，故 $K_a^q$ \Term{isomorphism}{同构}于单位 $\kappa(a)$-\Term{group}{群}，所以 $G_a$ \Term{solvable}{可解}（相应地，\Term{nilpotent}{幂零}）。
\end{proof}

\SGAsection{9}{商层}
本节主要在群的\Term{homogeneous-space}{齐性空间}这一特殊情形下，回顾关于按\Term{flat}{平坦}\Term{equivalence-relation}{等价关系}取商的一些熟知的一般事实（参见第 IV 讲）。

\begin{SGAresult}{定义}{9.1}
给定 $S$-\Term{group}{群}的一个\Term{monomorphism}{单态射} $u:G'\to G$，将 $G$ 关于下述\Term{monomorphism}{单态射}定义的\Term{equivalence-relation}{等价关系}的商\Term{sheaf}{层}（对于 (fpqc) \Term{topology}{拓扑}）记作 $G/G'$（相应地，$G'\backslash G$），并称为 $G$ 对 $G'$ 的\Term{right-quotient-sheaf}{右商层}（相应地，\Term{left-quotient-sheaf}{左商层}）：
\[
G\times_S G'\xrightarrow{\,\delta\circ(\operatorname{id}_G\times u)\,}G\times_S G,
\qquad
\left(\text{相应地， }G'\times_S G\xrightarrow{\,\gamma\circ(u\times\operatorname{id}_G)\,}G\times_S G\right),
\]
其中 $\delta$（相应地，$\gamma$）表示 $G\times_S G$ 的\Term{automorphism}{自同构}，它对 $g,h\in G(T)$ 由 $(g,h)\mapsto(g,gh)$（相应地，$(h,g)\mapsto(hg,g)$）定义。
\end{SGAresult}

\begin{SGAresult}{命题}{9.2}
设 $u:G'\to G$ 为一个 $S$-\Term{group}{群}\Term{monomorphism}{单态射}。假设 $G/G'$ 可由一个 $S$-\Term{scheme}{概形} $G''$ 表示。则：

(i) \Term{canonical}{典范}\Term{morphism}{态射} $p:G\to G''$ 对 (fpqc) \Term{topology}{拓扑}是\Term{covering-morphism}{覆盖态射}。

(ii) 若令 $\varepsilon''=p\circ\varepsilon$（此\Term{morphism}{态射}称为 $G''$ 的\Term{unit-section}{单位截面}），则以下图表都是\Term{cartesian-diagram}{笛卡儿图}：
\[
\begin{tikzcd}
G\times_S G'\arrow[r,"\mu\circ(\operatorname{id}_G\times u)"]\arrow[d,"\operatorname{pr}_1"']&G\arrow[d,"p"]\\
G\arrow[r,"p"']&G''
\end{tikzcd}
\qquad
\begin{tikzcd}
G'\arrow[r,"u"]\arrow[d,"\pi'"']&G\arrow[d,"p"]\\
S\arrow[r,"\varepsilon''"']&G''.
\end{tikzcd}
\]
特别地，$u$ 是一个\Term{immersion}{浸入}。

(iii) $G''$ 上存在唯一的左\Term{operator-group}{算子群}为 $G$ 的 $S$-\Term{scheme}{概形}结构，使 $p$ 为带\Term{operator-group}{算子群} $G$ 的 $S$-\Term{scheme}{概形}\Term{morphism}{态射}。

(iv) 若再假设 $G'$ 在 $G$ 中\Term{normal-subgroup}{正规}，则 $G''$ 上存在唯一的 $S$-\Term{group}{群}结构，使 $p$ 为 $S$-\Term{group}{群}\Term{morphism}{态射}。

(v) 设 $S_0$ 为一个 $S$-\Term{scheme}{概形}，令 $G_0=G\times_S S_0$、$G'_0=G'\times_S S_0$，则 $G_0/G'_0$ 可由 $G''_0=G''\times_S S_0$ 表示。\footnote{编者注（85）：也就是说，这个商是\Term{universal-quotient}{普遍商}，参见第 IV 讲第 3 节。}

(vi) 设 $\mathscr P$ 为 $S$-\Term{morphism}{态射}的一种性质。假设 $\mathscr P$ 在\Term{base-change}{基变换}下\Term{stable}{稳定}；若 $p:G\to G''$ 满足 $\mathscr P$，则\Term{structure-morphism}{结构态射} $\pi':G'\to S$ 也满足 $\mathscr P$。

\SourcePage{400}
(vii) 设 $\mathscr P$ 为 $S$-\Term{morphism}{态射}的一种性质。假设 $\mathscr P$ 对 (fpqc) \Term{topology}{拓扑}具有\Term{local-character}{局部性}（参见 2.0 与 2.1.2）。则\Term{morphism}{态射} $p:G\to G''$ 满足 $\mathscr P$，当且仅当 $\pi':G'\to S'$ 满足 $\mathscr P$。

(viii) 设 $\mathscr P$ 为 $S$-\Term{morphism}{态射}的一种性质；假设 $\mathscr P$ 对 (fpqc) \Term{topology}{拓扑}具有\Term{local-character}{局部性}且在复合下\Term{stable}{稳定}。若\Term{structure-morphism}{结构态射} $G''\to S$ 与 $G'\to S$ 满足 $\mathscr P$，则\Term{structure-morphism}{结构态射} $G\to S$ 也满足 $\mathscr P$。

(ix) 假设 $G$ \Term{reduced}{约化}，则 $G''$ \Term{reduced}{约化}。

(x) $G''$ 在 $S$ 上\Term{separated}{分离}的充要条件是 $u$（\Term{equivalence}{等价}地，$\varepsilon''$）为\Term{closed-immersion}{闭浸入}。

(xi) $G'$ 在 $S$ 上\Term{flat}{平坦}的充要条件是 $p$ 为\Term{flat}{平坦}\Term{morphism}{态射}（\Term{equivalence}{等价}地，\Term{faithfully-flat}{忠实平坦}）。

此时，$G''$ 在 $S$ 上\Term{flat}{平坦}的充要条件是 $G$ 在 $S$ 上\Term{flat}{平坦}。

(xii) $G'$ 在 $S$ 上\Term{flat}{平坦}且\Term{locally-finitely-presented}{局部有限表示}的充要条件是 $p:G\to G''$ \Term{faithfully-flat}{忠实平坦}且\Term{locally-finitely-presented}{局部有限表示}。

此时，要使 $G''$ 在 $S$ 上\Term{locally-finitely-presented}{局部有限表示}（相应地，\Term{locally-finite-type}{局部有限型}、\Term{finite-type}{有限型}、\Term{smooth}{光滑}、\Term{etale}{平展}、\Term{unramified}{非分歧}、\Term{locally-quasi-finite}{局部拟有限}、\Term{quasi-finite}{拟有限}），只须 $G$ 在 $S$ 上具有相同性质（前两种情形中，该条件也必要，参见 (viii)）。

(xiii) 假设 $G'$ 在 $S$ 上\Term{flat}{平坦}且\Term{finitely-presented}{有限表示}。

a) 则 $p$ \Term{finitely-presented}{有限表示}且\Term{faithfully-flat}{忠实平坦}。

b) 此外，$G$ 在 $S$ 上\Term{finitely-presented}{有限表示}，当且仅当 $G''$ 在 $S$ 上\Term{finitely-presented}{有限表示}。
\end{SGAresult}
\begin{proof}
回顾，所考虑的\Term{equivalence-relation}{等价关系}是\Term{universally-effective}{普遍有效}的（IV 4.4.9）。于是，断言 (i)、(iii)、(iv)、(v) 及 (ii) 的第一条断言分别由 IV 4.4.3、5.2.2、5.2.4、3.4.5 及 3.3.2 (iii) 得出。(ii) 的第二条断言由第一条得出，如下面的\Term{cartesian-diagram}{笛卡儿图}所示，因为 $(G\times_S G')\times_G S$ \Term{isomorphism}{同构}于 $G'$：
\[
\begin{tikzcd}
G'\arrow[r,"{((\varepsilon\circ\pi'),\operatorname{id}_{G'})}"]\arrow[d,"\pi'"']
&G\times_S G'\arrow[r,"\mu\circ(\operatorname{id}_G\times u)"]\arrow[d,"\operatorname{pr}_1"']
&G\arrow[d,"p"]\\
S\arrow[r,"\varepsilon"']&G\arrow[r,"p"']&G''.
\end{tikzcd}
\]
最后，显然 $\varepsilon''$ 是 $G''$ 的一个 $S$-\Term{section}{截面}，故为\Term{immersion}{浸入}（EGA I，5.3.13）；由上面的\Term{cartesian-diagram}{笛卡儿图}，$u$ 也如此，这就完成了 (ii) 的证明。此外，(vi) 立即由 (ii) 的第二个\Term{cartesian-diagram}{笛卡儿图}得出。

证明 (vii)。由 (i)，$p$ 对 (fpqc) \Term{topology}{拓扑}为\Term{covering-morphism}{覆盖态射}；所以由 (ii)，为证明 $p$ 满足 $\mathscr P$，只须证明第一\Term{projection}{投影} $\operatorname{pr}_1:G\times_S G'\to G$ 满足 $\mathscr P$。由于 $\mathscr P$ 在\Term{base-change}{基变换}下\Term{stable}{稳定}，而 $\operatorname{pr}_1$ 由 $\pi'$ 经\Term{base-change}{基变换}得到，这一点成立。

\SourcePage{401}
显然 (viii) 由 (vii) 得出，因为 $\pi=\pi''\circ p$，其中 $\pi'':G''\to S$ 表示\Term{structure-morphism}{结构态射}。

证明 (ix)。由 (i)，$p$ 是一个\Term{epimorphism}{满态射}；由于 $G$ \Term{reduced}{约化}，$p$ 经过\Term{immersion}{浸入} $G''_{\mathrm{r\text{é}d}}\to G''$ \Term{resolution}{分解}，所以后者也是\Term{epimorphism}{满态射}，因而为\Term{isomorphism}{同构}（IV 4.4.4）。

证明 (x)。若 $G''$ 在 $S$ 上\Term{separated}{分离}，则由 EGA I，5.4.6，$\varepsilon''$ 为\Term{closed-immersion}{闭浸入}。另一方面，(ii) 已表明，$\varepsilon''$ 为\Term{closed-immersion}{闭浸入}，当且仅当 $u$ 为\Term{closed-immersion}{闭浸入}。最后，若 $u$ 为\Term{closed-immersion}{闭浸入}，则 $\delta\circ(\operatorname{id}_G\times u):G\times_S G'\to G\times_S G$ 也如此；故由下面的\Term{lemma}{引理} 9.2.1，$G''$ 在 $S$ 上\Term{separated}{分离}。

断言 (xi) 由 (vii) 和 EGA $\mathrm{IV}_2$，2.2.13 得出。

断言 (xii) 由 (vii)、EGA $\mathrm{IV}_4$，17.7.5 与 17.7.7 及如下事实得出：由于 $p$ \Term{universally-open}{普遍开}，对任意 $s\in S$，若 $G_s$ 的底\Term{space}{空间}\Term{discrete}{离散}，则 $G''_s$ 的底\Term{space}{空间}也\Term{discrete}{离散}。

最后，断言 (xiii) 由 (vii)、(viii) 和 EGA $\mathrm{IV}_4$，17.7.5 得出。
\end{proof}

\begin{SGAresult}{\Term{lemma}{引理}}{9.2.1}
设 $X$ 为一个 $S$-\Term{scheme}{概形}，$R$ 为由\Term{monomorphism}{单态射} $v:R\to X\times_S X$ 在 $X$ 上定义的一个\Term{equivalence-relation}{等价关系}。假设 $R$ \Term{effective-equivalence-relation}{有效}。则：

(i) $v$ 是一个\Term{immersion}{浸入}；若 $Y=X/R$ \Term{separated}{分离}，则它是\Term{closed-immersion}{闭浸入}。

(ii) 再假设 $Y$ 表示 $X$ 对 $R$ 的 (fpqc) 商\Term{sheaf}{层}，\footnote{编者注（86）：如 O. Gabber 所指出的，证明中用到了这一点，必须将它加入假设。}且 $v$ 为\Term{closed-immersion}{闭浸入}。则 $Y=X/R$ 在 $S$ 上\Term{separated}{分离}。
\end{SGAresult}
\begin{proof}
回顾（IV，定义 3.3.2），“$R$ \Term{effective-equivalence-relation}{有效}”这一假设是指存在一个 $S$-\Term{scheme}{概形}\Term{morphism}{态射} $p:X\to Y$，使自然\Term{morphism}{态射} $R\to X\times_Y X$ 为\Term{isomorphism}{同构}。由此得到以下\Term{cartesian-diagram}{笛卡儿图}（EGA I，5.3.5）：
\[
\begin{tikzcd}
R\arrow[r,"v"]\arrow[d]&X\times_S X\arrow[d,"p\times p"]\\
Y\arrow[r,"\Delta_{Y/S}"']&Y\times_S Y.
\end{tikzcd}
\]
由于 $\Delta_{Y/S}$ 是\Term{immersion}{浸入}（EGA I，5.3.9），$v$ 也如此。同样，若 $Y$ 在 $S$ 上\Term{separated}{分离}，则 $\Delta_{Y/S}$ 为\Term{closed-immersion}{闭浸入}，$v$ 也如此。

反过来，假设 $v$ 为\Term{closed-immersion}{闭浸入}，且 $Y$ 表示 $X$ 对 $R$ 的 (fpqc) 商\Term{sheaf}{层}。则 $p$ 对 (fpqc) \Term{topology}{拓扑}为\Term{covering-morphism}{覆盖态射}（IV 4.4.3），故 $p\times p$ 也如此（由\Term{base-change}{基变换}，$p\times\operatorname{id}_X$ 与 $\operatorname{id}_Y\times p$ 为\Term{covering-morphism}{覆盖态射}，所以它们的复合 $p\times p$ 也如此）。因此，由 (fpqc) \Term{descent}{下降}（参见 EGA $\mathrm{IV}_2$，2.7.1），$\Delta_{Y/S}$ 是\Term{closed-immersion}{闭浸入}，即 $Y$ 在 $S$ 上\Term{separated}{分离}。
\end{proof}

\textbf{注记 9.2.2。}在 9.2 的一般假设下，若假设 $G'$ 在 $S$ 上\Term{flat}{平坦}且\Term{locally-finitely-presented}{局部有限表示}，则由 9.2 (vii)，$p$ 对 (fppf) \Term{topology}{拓扑}为\Term{covering-morphism}{覆盖态射}，
\SourcePage{402}
\footnote{编者注（87）：将 (fpqc) 更正为 (fppf)。}因此 9.2 的断言 (vii)、(viii) 可推广到对于 (fppf) \Term{topology}{拓扑}具有\Term{local-character}{局部性}的性质 $\mathscr P$。

\textbf{注记 9.3。}a) 商 $G/G'$ 是否\Term{representable}{可表}往往是一个微妙的问题；在本研讨班中，我们将证明某些特殊商的\Term{representability}{可表性}。

一般而言，为断言商 $G/G'$ \Term{representable}{可表}，仅假设 $G,G'$ 在 $S$ 上\Term{finitely-presented}{有限表示}且 $G'$ 在 $S$ 上\Term{flat}{平坦}是不够的。事实上，再假设 $G$ \Term{smooth}{光滑}且\Term{connected}{纤维连通}。此时，若 $G/G'$ 是一个\Term{scheme}{概形}，由推论 5.4，它\Term{separated}{分离}；由 9.2 (x)，$G'\hookrightarrow G$ 为\Term{closed-immersion}{闭浸入}。因此，若 $G'$ 在 $G$ 中不闭，则 $G/G'$ 不\Term{representable}{可表}。

为得到这样一个反例，可以取 $S$ 为一个\Term{discrete-valuation-ring}{离散赋值环}的谱，并令 $G=\mathbf G_{\mathrm m,S}$。再考虑一个在 $S$ 上\Term{invertible}{可逆}的整数 $n>1$，则 $\boldsymbol\mu_n=\operatorname{Ker}(G\xrightarrow{n}G)$ 是 $G$ 的一个在 $S$ 上\Term{etale}{平展}的\Term{closed-subgroup}{闭子群}（参见 VIIA）。\footnote{编者注（88）：参见 VIIA，8.4 或 VIII，2.1。}令 $G'$ 为 $\boldsymbol\mu_n$ 的如下\Term{open-subgroup}{开子群}：从 $\boldsymbol\mu_n$ 中删去其\Term{closed-fiber}{闭纤维}内原点的补集这一闭集。则 $G'$ 在 $G$ 中不闭，所以 $G/G'$ 不\Term{representable}{可表}。（也可以构造这样的例子，使 $G'$ \Term{smooth}{光滑}且\Term{connected}{纤维连通}。）

b) 反之，当 $G,G'$ 在 $S$ 上\Term{finitely-presented}{有限表示}，且 $G'$ 在 $S$ 上\Term{flat}{平坦}并在 $G$ 中闭时，或许 $G/G'$ 总是\Term{representable}{可表}。\footnote{这过于乐观，M. Raynaud 在其学位论文（同上，X 14）中已说明这一点。编者注（89）：此注指 [Ray70a] 中的反例 X.14，其中\Term{base}{基底}是\Term{dimension}{维数}为 $2$ 的\Term{regular}{正则}\Term{local-scheme}{局部概形}。}在这些假设下，已知下列特殊情形中 $G/G'$ \Term{representable}{可表}：

$1^\circ$——$S$ 为一个 Artin 环的谱（参见 VIA 3.2 和 3.3.2）。

$2^\circ$——$G'$ 在 $S$ 上\Term{proper}{固有}，且 $G$ 在 $S$ 上\Term{quasi-projective}{拟射影}（参见 V 7.1）。

$3^\circ$——$S$ \Term{locally-noetherian}{局部 Noether}且\Term{dimension}{维数}为 $1$（参见 [An73]，\Term{theorem}{定理} 4.C）。
